% SPDX-FileCopyrightText: 2026 Will Cook
% SPDX-License-Identifier: CC-BY-4.0
%
% Standalone Erdős Problem Note for #1041.  Context is deliberately repeated
% from the sibling notes so that this file remains independently readable.
%
% The public March manuscript is attributable only to the erdosproblems.com
% handle `shtuka`; the PDF contains no real-name author metadata.  The
% bibliography therefore cites that handle and public URL conservatively.
%
% Build from the public paper directory with `make`.
\documentclass[11pt]{article}
\usepackage{longtable}
% SPDX-FileCopyrightText: 2026 Will Cook
% SPDX-License-Identifier: CC-BY-4.0
%
% Shared preamble for the Erdős Problem Notes: the short, standalone,
% problem-owned manuscripts covering the expansion library `ErdosProblems`.
%
% One note owns one Erdős problem.  Each note repeats whatever context its own
% reader needs, so that a reader who arrives at a single problem never has to
% assemble it from the gateway paper.  Deliberate repetition across notes is the
% design, not an oversight.
%
% Source links resolve against one pinned formal-source commit, declared once
% below.  `scripts/check_problem_note_sources.py` verifies that every authored
% (file, line, declaration) triple names that exact declaration in the pinned
% snapshot, so a link cannot silently rot as later work moves lines.

\usepackage{paper-house-style}
\usepackage{array}
\usepackage{booktabs}
\usepackage{enumitem}
\setlist{topsep=0.35em,itemsep=0.2em}
\usepackage[colorlinks=true,linkcolor=linkinternal,urlcolor=linkexternal,
  citecolor=linkinternal]{hyperref}
\hypersetup{bookmarksnumbered=true,bookmarksopen=true,bookmarksopenlevel=1,
  pdfauthor={Will Cook},pdflang={en-GB}}

% ---------------------------------------------------------------------------
% Pinned formal source.  \commit names the checkpoint every link resolves
% against; the checker reads that snapshot rather than the working tree, so the
% printed coordinates stay correct however later waves move declarations.
% ---------------------------------------------------------------------------
\newcommand{\commit}{e5fd26a6d1b1a346430205eab5385b4c9565d004}
\newcommand{\commitshort}{e5fd26a6d1b1}
\newcommand{\repobase}{https://github.com/wcook04/plectis-lean-erdos249-257/blob/\commit}
\newcommand{\sourceurl}{https://github.com/wcook04/plectis-lean-erdos249-257/tree/\commit}
\newcommand{\PX}{ErdosProblems}

% Declaration names stay inside link targets.  The printed label is either a
% spaced, lower-case reading of the declaration name (\lref) or a short authored
% phrase (\lword); neither prints a module path or a raw Lean identifier.
\ExplSyntaxOn
\cs_new_protected:Npn \noteleanlabel #1
  {
    \tl_set:Nx \l_tmpa_tl { \tl_to_str:n {#1} }
    \regex_replace_all:nnN { _+ } { \c{space} } \l_tmpa_tl
    \regex_replace_all:nnN
      { ([a-z0-9])([A-Z]) }
      { \1\c{space}\2 }
      \l_tmpa_tl
    \tl_set:Nx \l_tmpa_tl { \tl_lower_case:n { \tl_use:N \l_tmpa_tl } }
    \tl_use:N \l_tmpa_tl
  }
\ExplSyntaxOff
\newcommand{\leanlabel}[1]{\noteleanlabel{#1}}
\newcommand{\lref}[3]{\href{\repobase/\PX/#1\#L#2}{\leanlabel{#3}}}
\newcommand{\lword}[4]{\href{\repobase/\PX/#1\#L#2}{#4}}
\newcommand{\lloc}[2]{\href{\repobase/\PX/#1\#L#2}{Lean source}}

% Links into a sibling library, named repository-relative.  The expansion
% library is implicit in \lref/\lword, because that is where problem-owned
% work lands.  Several problems have their headline theorem in the older
% reviewed #249/#257 corpus, and a note that could not name that library
% would have to paraphrase a checked statement instead of linking it.
\newcommand{\mref}[3]{\href{\repobase/#1\#L#2}{\leanlabel{#3}}}
\newcommand{\mword}[4]{\href{\repobase/#1\#L#2}{#4}}
\newcommand{\mloc}[2]{\href{\repobase/#1\#L#2}{Lean source}}
\emergencystretch=2em

\newtheorem{theorem}{Theorem}[section]
\newtheorem{proposition}[theorem]{Proposition}
\newtheorem{lemma}[theorem]{Lemma}
\newtheorem{corollary}[theorem]{Corollary}
\theoremstyle{definition}
\newtheorem{definition}[theorem]{Definition}
\newtheorem{example}[theorem]{Example}
\newtheorem{problem}[theorem]{Problem}
\theoremstyle{remark}
\newtheorem*{remark}{Remark}

\newcommand{\N}{\mathbb{N}}
\newcommand{\Npos}{\mathbb{N}_{>0}}
\newcommand{\Z}{\mathbb{Z}}
\newcommand{\Q}{\mathbb{Q}}
\newcommand{\R}{\mathbb{R}}
\newcommand{\lcm}{\operatorname{lcm}}
\newcommand{\ph}{\varphi}

% The series line printed under every note's title block.
\newcommand{\seriesline}{Erd\H{o}s Problem Notes}

% Production provenance.  The wording is fixed by the corpus provenance
% contract and reproduced verbatim in every manuscript in this repository, so
% the papers cannot drift into different accounts of how they were made.
\newcommand{\noteprovenance}{\provenancenote{\emph{Authorship and AI use.} The
prose, formal proofs, and software in this version were drafted and revised by
large-language-model agents under Will Cook's direction.  Cook set the
objectives and acceptance criteria, reviewed and authorised the manuscript, and
is responsible for its contents.  The agents are production tools, not authors.
See \emph{Statements and declarations}.}}

% The status paragraph every note carries, in its own words, before any result.
% Kernel checking establishes the exact Lean proposition; it does not establish
% that the proposition is the interesting one, that it is new, or that the
% Erdős problem is closed.
\newcommand{\statusboundary}{%
  \emph{Status.}  The problem treated here is open, and this note does not
  close it.  Every statement below marked as checked is a proposition that the
  pinned Lean kernel accepts from the sources this note links to, with no
  \texttt{sorry}, no added axiom, and no unchecked evaluation.  That is a claim
  about the formal statement, not about its mathematical interest, its
  novelty, or the original problem.  The unresolved obligations are named
  exactly, in their own section, and none of the finite computations,
  reductions, or no-go results here removes one of them.}

\hypersetup{
  pdftitle={Sharp solved families and constant-factor paths in polynomial lemniscates for Erdős Problem 1041},
  pdfsubject={Erdős Problem Notes: Problem 1041},
  pdfkeywords={lemniscate, critical values, Newton flow, free-point inequality, saddle connection, Lean 4},
}

% The prose keeps compact declaration labels; the source appendix supplies the
% commit-pinned links checked by the public release gate.
\newcommand{\pdecl}[1]{\emph{\leanlabel{#1}}}
\newcommand{\dbar}{\mathbb{D}}

\runninghead{Erd\H{o}s \#1041: solved families and constant-factor paths}
\title{Sharp Solved Families and\\Constant-Factor Paths in Polynomial Lemniscates}
\subtitle{Unconditional parent progress, exact obstructions, and frontier
mechanisms for Erd\H{o}s Problem~\#1041}
\author{Will Cook}
\date{30 August 2026}

\begin{document}
\maketitle
\noteprovenance

\begin{abstract}
\noindent
Let $f(z)=\prod_i(z-z_i)$ be monic, with every root in the open unit disc.
Erd\H{o}s, Herzog, and Piranian asked whether two roots can always be joined
inside $\{|f|<1\}$ by a curve of length less than $2$.  The problem remains
open.  Our strongest theorem applying to arbitrary degree and arbitrary root
geometry is an unconditional constant-factor version.  Put
\[
 \mu=\min_{f'(c)=0}|f(c)|,\qquad \rho=\mu^{1/n}.
\]
Two zero occurrences are joined inside $\{|f|\le2\mu\}$ by a path of length
at most $(71/10)\rho$.  If the roots lie in the unit disc and $\mu\le1/2$,
the construction lies in $\{|f|<1\}$ and has length at most $5.7$.
It proves the target constant $2$ whenever the first-merge component has at
least $17$ roots, and also under explicit arity-dependent component-capacity
defect inequalities.  These are ordinary analytic theorems, not Lean
formalisations; no literature-priority claim is made.

We also prove the conjectured path bound for three classes of polynomials.
For collinear roots the result is all-degree and sharp: one adjacent-root
segment works with the best Chebyshev constant.  The other complete classes
are translated cubic quotient fibres $P((z-h)^q)$ and primitive sparse
quintics $z^5+az^4+bz+c$.  Lean checks their principal selectors and extremal
inequalities; the finite-fibre, path-assembly, and affine-transport arguments
are proved in the paper.

At every non-root critical point $c$ we prove the sharp metric selector: if
$n\ge2$ and
\[
 r^n=\prod_{k=1}^n|c-z_k|,
\]
then two distinct roots satisfy
\[
 |c-z_i|+|c-z_j|\le2r.
\]
Lean checks this theorem.  It also checks two counterexamples showing why the
bound does not immediately give the required path: the nearest critical spoke
can leave the unit lemniscate, and some cubics admit no valid pair of straight
spokes through the critical point.

Finally, we identify two specific obstacles to the remaining argument.  The
merge-tree lifetime is governed by the strictly convex function
$\Phi(x)=\int_0^x(\log(\coth t))^{-1}\,dt$, for which
$\Phi(x)/x\to0$ as $x\to0$; no positive uniform linear conversion from
attachment age to lifetime is possible.  Separately, a printed spanning-tree
argument fails both at a four-pronged saddle and on an exact Cassini example.
The unresolved work is therefore to control the relevant merge ratios and to
construct a compatible planar path with total length below $2$.
\end{abstract}

\section{The problem}\label{sec:problem}

\begin{problem}[Erd\H{o}s \#1041]\label{res:problem}
Let $f(z)=\prod_{i=1}^{n}(z-z_i)$ be monic with $z_i\in\dbar$, the open unit
disc.  Show that two of the roots can be joined by a curve of length less than
$2$ lying in the open lemniscate $E=\{z\in\mathbb{C}:|f(z)|<1\}$.
\end{problem}

\noindent Numbering and current status follow Bloom's Erd\H{o}s problem
catalogue \cite{bloom}.
The problem is open.  The original source is Problem~5 on printed p.~139 of
Erd\H{o}s--Herzog--Piranian \cite[p.~139]{ehp1958}; the preceding paragraph records the
known input that one component of the lemniscate contains at least two zeros.

Two recent manuscripts are relevant.  The 48-page manuscript posted by
\texttt{shtuka} on 24 March 2026
\cite[Theorem~1, p.~1]{march2026} claimed the unrestricted
statement.  Its Proposition~12 (p.~16, with proof continuing through p.~30)
supplies the spanning-tree decomposition used in its final proof.  That
manuscript has been withdrawn by its author, who wrote on the problem's
discussion page on 26 March 2026 that the statement of Proposition~12, and not
only its proof, is incorrect, after Tao identified an unjustified invocation of
its Lemma~8 the previous day.  The file has not been revised since 24 March
2026.  Section~\ref{sec:gap} gives an exact Cassini counterexample to
Proposition~12's metric conclusion, reached independently; the same witness
refutes Proposition~7 of an unattributed local revision, and the local
three-ended saddle construction fails separately.  A second general-degree
attempt circulated on 22 April 2026 and was broken on 27 April 2026.  No
general-degree manuscript stands, and the problem page records the problem as
open with no accepted proof claim.  Pendyala's independent June 2026
preprint \cite[Thm.~1, p.~1]{june2026} proves the degree-four case.  The quartic theorem does
not close the problem.  An adjacent June 2026 paper by Pendyala
\cite[Thm.~1.2, p.~2; introductory discussion, pp.~1--2]{pendyala2026shortest}
proves $c\sqrt{\log n}\leq S(n)\leq\pi n$ for all sufficiently large $n$,
where $S(n)$ measures the shortest path from $0$ to $\partial\mathbb D$ in
the normalized set $\{|z|\leq1,\ |f(z)|\leq1\}$.  This is a different
statement from the root-to-root problem here: it supplies important adjacent
lemniscate-path context, but it is not evidence that Erd\H{o}s~\#1041 is
solved or even a partial result for its root-pair conclusion.

The two primary comparisons are pinned to downloaded PDFs.  The
Erd\H{o}s--Herzog--Piranian source and digest are
\begin{center}
\footnotesize
\path{annexes/erdos-herzog-piranian-1958-metric-properties-of-polynomials/source.pdf}\\
SHA-256: \nolinkurl{bff39876de5e152b5a0d2b622eb6f9ca0c747d401ceea68146cfa2b2b5f28ce9}.
\end{center}
Its exact locator is printed p.~139 / PDF p.~15 / extracted lines 594--602,
where the unit-disc setting, the two-zero component, and Problem~5 appear.
It supports the historical problem statement only, not the local length--area
estimate or the all-degree constant-factor result.  Pendyala's downloaded
degree-four source and digest are
\begin{center}
\footnotesize
\path{annexes/arxiv-2606-24875-pendyala-degree-four-lemniscate-path/source.pdf}\\
SHA-256: \nolinkurl{337db0cc046023b1429d9131720d199843766629cd2ad8bf8b04a40168d7bcd8}.
\end{center}
Its abstract and Theorem~1 are on PDF p.~1, the four-point radial lemma is on
PDF pp.~1--2, and the theorem proof runs through p.~3.  That proof's
minimum-enclosing-disk/four-point mechanism is a genuine degree-four result,
but it does not extend the all-degree claims or validate the Newton-flow
decomposition below.

The adjacent shortest-path comparison is likewise pinned to its downloaded
primary source.  Pendyala's 34-page arXiv:2606.19178v1 PDF was downloaded from
\url{https://arxiv.org/pdf/2606.19178v1}; the arXiv API record is
\url{https://export.arxiv.org/api/query?id_list=2606.19178}.  Its digest is
\begin{center}
\footnotesize
SHA-256: \nolinkurl{d7902a8bc37d2aad6dea52dffbbbaaf7ff23954ece5a0c4cbab108f58df4e6df}.
\end{center}
The abstract (printed p.~1, extracted lines 12--32) and Theorem~1.2 (printed
pp.~3--4, extracted lines 97--113) support the bounds
$c\sqrt{\log n}\leq S(n)\leq\pi n$.  The introduction's distinction from the
root-to-root problem is on printed pp.~2--3 (extracted lines 51--78); the
reciprocal-sweep and alternating-maze mechanisms are described in extracted
lines 79--96.  These locators support only the neighboring marked-point result:
they do not support a root-to-root path, a degree-free bound, or a transfer to
the first-merge problem.

The object we work with is not the lemniscate directly but the flow that
foliates it.

\statusboundary

\subsection*{Main results at a glance}

The problem remains open.  The table is ordered by mathematical consequence:
sharp solved families and unconditional progress on the parent geometry come
before auxiliary selectors and mechanism eliminations.
Section~\ref{sec:evidence-ledger} gives the complete evidence and boundary
ledger.

{\small
\renewcommand{\arraystretch}{1.18}
\begin{tabular}{@{}
  >{\raggedright\arraybackslash}p{0.29\linewidth}
  >{\raggedright\arraybackslash}p{0.63\linewidth}@{}}
\toprule
\textbf{Result} & \textbf{What is established} \\
\midrule
\textbf{Sharp collinear family} &
The conjecture holds in every degree for collinear roots, with the sharp
Chebyshev constant.  \textsc{Checked kernels; ordinary affine transport.} \\
\addlinespace
\textbf{Global constant-factor theorem} &
For arbitrary monic $f$, two zero occurrences are joined in
$\{|f|\le2\mu\}$ by a path of length at most
$(71/10)\mu^{1/n}$.  For unit-disc roots and $\mu\le1/2$, the path lies in
$\{|f|<1\}$ and has length at most $5.7$.
\textsc{Ordinary exact proof.} \\
\addlinespace
\textbf{Parent-theorem regimes} &
The target constant $2$ holds when $\mu\le1/2$ and the first-merge component
has at least $17$ roots, or under an explicit arity-dependent
component-capacity defect bound.  \textsc{Ordinary exact proof.} \\
\addlinespace
\textbf{Critical-value separation regime} &
If a simple critical point $c$ with $v=f(c)\ne0$ has every other critical
point $d$ at normalised distance $|1-f(d)/v|\ge S>1$, two roots are joined
inside $\{|f|\le|v|\}$ by a path of length below $2|v|^{1/n}$ whenever
$(1+S)^{2/n}\log(S/(S-1))<1$.  The exact regimes are $S=4$ for $n\ge3$,
$S=3$ for $n\ge4$, and $S=2$ for $n\ge6$, and each cutoff is the first degree
at which its inequality holds.  \textsc{Ordinary exact proof; Lean checks the
degree monotonicity, the three thresholds, and the squared-length consumer.} \\
\addlinespace
\textbf{Critical-spectrum closure} &
At a simple minimising critical point, enough separation from the next
critical-value modulus gives a join inside the sharper level $K_\mu$ and an
explicit degree-dependent criterion for the target constant $2$.
\textsc{Ordinary exact proof.} \\
\addlinespace
\textbf{Primitive sparse quintics} &
Two exact tail-energy selectors give the conjectured path for the open-disc
sparse-quintic family.  \textsc{Checked selectors; ordinary path assembly.} \\
\addlinespace
\textbf{Translated cubic quotient fibres} &
The conjectured path exists for every $f(z)=P((z-h)^q)$ with $P$ monic cubic
and $q\ge2$.  \textsc{Checked kernel; ordinary finite-fibre assembly.} \\
\addlinespace
\textbf{All-degree proximity theorem} &
At every non-root critical point, two roots have total critical-point
distance at most twice the geometric mean of all root distances.
\textsc{Lean checked.} \\
\addlinespace
\textbf{Straight-path obstructions} &
Exact quintic and cubic families rule out the unique-nearest spoke and every
pair of straight spokes as universal constructions.  \textsc{Lean checked.} \\
\addlinespace
\textbf{Current boundary} &
The general problem is still open.  A counterexample in the
$\mu\le1/2$ branch must have first-merge arity at most $16$ and lie above the
corresponding capacity-defect cutoff; the high-critical branch also remains. \\
\bottomrule
\end{tabular}}

\section{An unconditional all-degree constant-factor path theorem}
\label{sec:constant-factor}

The parent problem asks simultaneously for containment in the open unit
sublevel and the sharp length $2$.  The following theorem keeps the same
root-to-root path object for arbitrary geometry and arbitrary degree, while
relaxing both constants in a controlled scale-invariant way.  For a monic
degree-$n$ polynomial put
\[
 K_t=\{z:|f(z)|\le t\},\qquad
 \mu=\min_{f'(c)=0}|f(c)|,\qquad \rho=\mu^{1/n}.
\]

\begin{theorem}[unconditional constant-factor path]
\label{res:constant-factor-path}
For every monic polynomial $f$ of degree $n\ge2$, two zero occurrences are
joined by a possibly degenerate path of length at most
\[
 \frac{71}{10}\,\rho
\]
inside $K_{2\mu}$.  If $f$ is squarefree, their locations are distinct.
If every root lies in the open unit disc and $\mu\le1/2$, the construction
may be chosen inside $\{|f|<1\}$ with length at most $5.7$.
\end{theorem}

The constant is the rationally certified specialization of a two-parameter
bound.  If the selected working component contains $k\ge2$ roots, then for
every $r\in(0,1)$ and $\lambda>1$ the proof constructs a path in
$K_{\lambda\mu}$ whose length is at most
\[
 \sqrt{\frac2k}\left(
   \frac{\sqrt2\,r}{(1-r)^2}
   +\lambda^{1/n}\left(
      \sqrt{\log(\lambda/r)}+
      \frac{\pi}{\sqrt{\log\lambda}}
    \right)
 \right)\rho.                                      \tag{CF}
\]
For $n\ge3$, the choice $\lambda=2$, $r=3/20$ makes the bracket less
than $71/10$; degree two has the exact root-segment bound $2\rho$.

\begin{proof}[Proof architecture]
Choose a critical point at level $\mu$ and a component at a regular level
$t\in[\mu,\lambda\mu]$ that contains the corresponding first merge.  If a
regular component $C'$ at level $\sigma$ contains $k'$ roots, winding and
Cauchy--Schwarz give the local perimeter estimate
\[
 \mathcal H^1(\partial C')^2
 \le2\pi k'\,\sigma\,\frac{d}{d\sigma}
      \operatorname{Area}(K_\sigma\cap C').         \tag{CF1}
\]
P\'olya's area--capacity inequality
$\operatorname{Area}(K_T)\le\pi T^{2/n}$ then selects $t$ with a
degree-free boundary-length bound.

For an argument avoiding the finitely many critical-value rays, lift the
radial value segment from each root to $\partial C_t$.  Split each lift at
level $r\mu$.  Koebe distortion controls the low pieces.  The exact
precritical conformal-radius energy inequality
\[
 \sum_i\frac{\mu^2}{|f'(z_i)|^2}\le\rho^2          \tag{CF2}
\]
aggregates those pieces without paying once per root.  Coarea and~(CF1)
control the high pieces on average over the argument.  Choose one argument
realising that average, order the lift endpoints cyclically on
$\partial C_t$, and use the shortest adjacent boundary arc.  Averaging the
two lifts and the boundary arcs over the $k$ adjacent pairs gives~(CF).
The strict unit-sublevel statement chooses the regular level below the top of
the window when $\lambda\mu\le1$.
\end{proof}

The same estimate reaches the sharp target in explicit regions of the parent
problem.  Let $k_0$ be the number of roots in the component of $K_\mu$
containing the chosen minimising critical point.  Let $C$ be the corresponding
component of $K_{2\mu}$ and put
\[
 \kappa=\frac{\operatorname{cap}(\overline C)}{(2\mu)^{1/n}}.
\]
Thus $\kappa=e^{-\Sigma/n}$, where $\Sigma$ is the total exterior Green
defect of the roots excluded from $C$.

\begin{theorem}[first-merge arity and capacity closures]
\label{res:first-merge-capacity-closures}
Assume all roots lie in the open unit disc and $\mu\le1/2$.

\begin{enumerate}[label=\textup{(\roman*)},leftmargin=*]
\item If $k_0\ge17$, two distinct roots are joined inside $\{|f|<1\}$ by a
path of length strictly less than $2$.
\item Put
\[
 A=\frac{283}{3610},\qquad B=\frac{52029}{9100},\qquad
 \tau_k=\frac{\sqrt{2k}-A}{B}.
\]
If $\kappa\le\tau_{k_0}$, the same conclusion holds.  In particular,
$\kappa\le1/3$ suffices for every $k_0\ge2$, while at $k_0=16$ the rational
cutoff improves to $\kappa\le39/40$.
\end{enumerate}
\end{theorem}

The second part retains the component root count and replaces the global area
bound in~(CF1) by
$\operatorname{Area}(C)\le\pi\operatorname{cap}(\overline C)^2$.
Consequently a counterexample with $\mu\le1/2$ must have $k_0\le16$ and must
lie above its arity-dependent capacity cutoff.  Smaller critical levels give
stronger nested regions: $\mu\le1/4$ closes every $k_0\ge12$,
$\mu\le1/8$ closes every $k_0\ge10$, and the $\mu\le1/64$ calculation raises
the surviving $k_0=9$ capacity cutoff to $99/100$.

There is a complementary closure criterion that spends separation in the
critical spectrum instead of component capacity.  Suppose $f$ is squarefree
and a critical point $c_*$ attaining $\mu$ is simple.  Let $T_*>\mu$ be the
supremum of the levels for which the component born at $c_*$ contains exactly
two roots.  Component root counts change only at critical-value moduli, so
$T_*$ is at least the next distinct critical-value modulus.

\begin{theorem}[minimal-hub spectral window]
\label{res:minimal-hub-window}
Under these hypotheses, two roots are joined by a path inside $K_\mu$ of
length at most
\[
 \Psi_n(T_*/\mu)\,\mu^{1/n},\qquad
 \Psi_n(x)=x^{1/n}
   \sqrt{2\log\frac{x+1}{x-1}}.                     \tag{SW}
\]
The function $\Psi_n$ is strictly decreasing on $(1,\infty)$.  Let
$x_n^*$ be the unique solution of $\Psi_n(x)=2$ for $n\ge3$.  If the roots
lie in a disk of radius $R$ and the second distinct critical-value modulus
$\nu_2$ satisfies $\nu_2/\mu>x_n^*$, then the path has length less than
$2R$ and proves the parent conclusion.  The first thresholds are
\[
 \begin{aligned}
 x_3^*&=1.6064943\ldots,&
 x_4^*&=1.4784372\ldots,\\
 x_5^*&=1.4289343\ldots,&
 x_{10}^*&=1.3598081\ldots .
 \end{aligned}
\]
and $x_n^*$ decreases to $1$.
\end{theorem}

Indeed, the two-root descent-arc estimate at any $t\in(\mu,T_*)$ is
\[
 \operatorname{len}(\Gamma)
 \le t^{1/n}\sqrt{2\log\frac{t+\mu}{t-\mu}},
 \qquad \Gamma\subset K_\mu.
\]
Writing $x=\coth u$ gives
$\Psi_n(x)^2=4u\coth(u)^{2/n}$ and proves strict decrease.  Taking the limit
$t\uparrow T_*$ gives~(SW); the Fekete--Vandermonde bound
$\mu\le R^n$ converts $\Psi_n(\nu_2/\mu)<2$ into the parent metric bound.
The same criterion is stable under the compositional pullback
$P((z-h)^q)$ after adjoining the hub value $|P(0)|$ to the critical spectrum.
The complete proof and threshold certificates are in
\path{ErdosProblems/Erdos1041/MinimalHubWindowJoin.md}.  This theorem is
silent at a multiple minimising critical point and makes no priority claim.

These statements are complete ordinary analytic proofs.  They are not Lean
theorems and make no literature-priority claim.  Their full constant
certificates, component-capacity tables, and local proofs are in
\path{ErdosProblems/Erdos1041/UnconditionalConstantFactorBound.md}; the
P\'olya input is the global area--capacity inequality, not the local
winding--coarea estimate, reciprocal-derivative aggregation, or assembled
path theorem.  The unrestricted constant $2$ and the high-critical or
low-arity/high-capacity residual cells remain open.

\section{Critical proximity and straight-path obstructions}
\label{sec:critical-proximity}

The strongest source-current theorem that applies in every degree is a sharp
metric selection principle at a critical point.  It consumes the complete
logarithmic-derivative balance rather than a degree-specific coefficient
pattern.

\begin{theorem}[critical geometric-mean proximity]
\label{res:critical-proximity}
Let $n\ge2$, let $z_1,\ldots,z_n,c\in\mathbb C$ with $c\ne z_k$ for every
$k$, and suppose
\[
  \sum_{k=1}^n\frac1{c-z_k}=0.
\]
If $r>0$ is determined by
\[
  r^n=\prod_{k=1}^n|c-z_k|,
\]
then there are distinct indices $i,j$ such that
\[
  |c-z_i|+|c-z_j|\le2r.
\]
\end{theorem}

\noindent The proof chooses the two smallest distances $d_i\le d_j$.
The reciprocal balance gives $d_j\le(n-1)d_i$, while minimality and the
geometric-mean identity give
$d_i d_j^{\,n-1}\le r^n$.  After normalising by $r$, a sharp two-variable
Bernoulli inequality yields $d_i+d_j\le2r$.  The complete complex theorem,
including the selector and both inequalities, is checked as
\pdecl{exists_two_roots_dist_sum_le_two_mul_geomMean}.

This is not silently promoted to a path theorem.  The same source module
checks two exact barriers to the most tempting completions.

\begin{proposition}[two straight-path no-go families]
\label{res:straight-no-go}
There is a monic quintic with all roots in the open unit disc and a non-root
critical point $c$ whose unique nearest root has a point on the straight spoke
to $c$ outside $\{|f|<1\}$.  There is also a monic cubic with all roots in the
open unit disc such that the midpoint of every pair of distinct roots lies
outside $\{|f|<1\}$.
\end{proposition}

\noindent The certificates are exact rational-algebraic computations, checked
as \pdecl{nearestSpoke_unique_nearest_spoke_escapes} and
\pdecl{allStraightCubic_every_pair_midpoint_escapes}.  The first defeats a
tie-breaking escape hatch; the second defeats every straight root-pair
segment.  Together with Theorem~\ref{res:critical-proximity}, they isolate the
real gap: critical-point proximity selects a short metric pair, but universal
containment requires curved or topological geometry.

\section{Solved polynomial families}
\label{sec:solved-families}

Three source-current families do cross that containment gap.  Their complete
paper theorems are ordinary mathematics assembled from formalised
load-bearing kernels; this section keeps that authority split explicit.

\subsection{Translated cubic quotient fibres}

\begin{theorem}[translated cubic quotient fibres]
\label{res:cubic-fibres}
Let $q\ge2$, $h\in\mathbb C$, and let $P$ be a monic cubic.  Put
\[
  f(z)=P((z-h)^q).
\]
If all zeros of $f$ lie in the open unit disc and $f$ has at least two
distinct zeros, then two distinct zeros of $f$ are joined by a two-segment
path of length strictly less than $2$ contained in $\{|f|<1\}$.
\end{theorem}

\noindent For cubic roots $r,s,v$, the real charges
$\Re(r\overline{s+v})$ sum to more than $-3$, so one is greater than $-1$.
For that root, AM--GM and the exact cyclotomic cancellation
$(1-t)(1+t+t^2)=1-t^3$ prove that its complete origin spoke is safe.  Lean
checks this fan-in as \pdecl{cubic_has_safe_root_spoke}.  Pulling the spoke
back under $y\mapsto y^q$ gives $q$ safe spokes from $h$; two different fibre
points are joined through $h$ with length $2|y|<2$.  The finite root-of-unity
mean-square argument and fibre assembly are ordinary, not hidden Lean claims.

\subsection{Primitive sparse quintics}

\begin{theorem}[primitive sparse quintics]
\label{res:primitive-quintic}
Let
\[
  p(z)=z^5+az^4+bz+c
\]
be monic with all five root occurrences in the open unit disc.  Then two
distinct root occurrences $w_i,w_j$ satisfy
\[
  |bw_i+c|<1,
  \qquad
  |bw_j+c|<1,
\]
and the corresponding two radial spokes join them through the origin inside
$\{|p|<1\}$ with total length strictly below $2$.  If the two occurrences
coincide, the associated root-to-root path is degenerate.
\end{theorem}

\noindent The stronger closed-disc selector has non-strict tails; if $a\ne0$
the selected tails are strict, while for $a=0$ strictness is exactly the
interior-root condition.  The phase-sensitive boundary theorem uses the first
three Newton moments and a cubic separator.  Its harmonic extension controls
mixed interior/boundary configurations and shows that an unsafe interior tail
contributes at most $2/31$.  Lean checks the finite selectors
\pdecl{primitiveBoundary_exists_two_tailEnergy_lt_one} and
\pdecl{primitiveInterior_exists_two_tailEnergy_lt_one}; the rotation,
moment identification, Abel path estimate, and two-spoke assembly are ordinary
mathematics.

\subsection{Sharp collinear roots}\label{bdry:solved-polynomial-families}

The complete all-degree collinear theorem appears in
Section~\ref{sec:collinear}.  It gives the sharp Chebyshev constant, not merely
an existence bound, and therefore belongs to the solved-family layer rather
than the frontier layer.  Lean checks the constrained alternation engine
\pdecl{exists_peak_le_of_monic_comparison} and its Chebyshev consumer
\pdecl{exists_peak_le_comparisonBound}; affine normalisation, root-gap
selection, and transport back to the original line are ordinary.

\section{Critical-value separation in every degree}
\label{sec:separation}

The solved families of Section~\ref{sec:solved-families} are unconditional and
complete, and each is confined to a shape of root configuration.  This section
gives a hypothesis that is not a shape at all.  It constrains only the critical
spectrum, it holds in every degree above an exact cutoff, and where it holds it
gives the target constant~$2$ rather than a constant multiple of it.

\begin{theorem}[critical-value separation]\label{res:critical-value-separation}
Let $f$ be monic of degree $n$, let $c$ be a simple critical point with
$v=f(c)\ne0$, and let $S>1$.  Suppose every other critical point $d$ of $f$
satisfies
\[
  \Bigl|1-\frac{f(d)}{v}\Bigr|\ge S .
\]
Then the square-root-resolved inverse branch through $c$ joins two roots of
$f$ inside $\{|f|\le|v|\}$ by a path of length at most
\[
  2\,|v|^{1/n}(1+S)^{1/n}\sqrt{\log\frac{S}{S-1}} .
\]
In particular the connector is shorter than $2|v|^{1/n}$ whenever
\begin{equation}\label{eq:separation-threshold}
  (1+S)^{2/n}\log\frac{S}{S-1}<1 .
\end{equation}
\end{theorem}

\begin{corollary}[the exact convenient regimes]\label{res:critical-value-thresholds}
Inequality~\eqref{eq:separation-threshold} holds for $S=4$ in every degree
$n\ge3$, for $S=3$ in every degree $n\ge4$, and for $S=2$ in every degree
$n\ge6$.  Each cutoff is the first degree at which its inequality holds: the
squared coefficient of \eqref{eq:separation-threshold} is $1.43841$ at
$(n,S)=(2,4)$, $1.02171$ at $(3,3)$, and $1.07566$ at $(5,2)$, while at
$(6,2)$ it is $0.99969$.
\end{corollary}

\begin{corollary}[a parent-problem regime]\label{res:separation-parent}
Let $f$ be monic with all roots in the open unit disc, let $c$ be a simple
critical point with $v=f(c)\ne0$, and suppose $|1-f(d)/v|\ge4$ for every other
critical point $d$.  If $n\ge3$ then Erd\H{o}s Problem~\#1041 holds for $f$.
\end{corollary}

\noindent The separation hypothesis makes $c$ the unique critical point of
minimum modulus, the Fekete chain gives $|v|<1$, and
Theorem~\ref{res:critical-value-separation} with $S=4$ places a connector of
length below $2$ inside $\{|f|<1\}$.

The proof of Theorem~\ref{res:critical-value-separation} is the same
area--capacity argument used in Section~\ref{sec:constant-factor}, run at a
single saddle instead of over a merge tree.  Normalise so that $c=0$, $v=1$,
and $|f''(0)|\ne0$.  The square substitution resolves the saddle, so
$f(Z(\xi))=1-\xi^2$ has two local holomorphic solutions interchanged by
$\xi\mapsto-\xi$.  A finite branch point of the algebraic continuation of $Z$
requires $f'(d)=0$ and $\xi^2=1-f(d)$, and the separation hypothesis puts every
such point on or outside $|\xi|=\sqrt S$; properness of a polynomial rules out
escape to infinity over a bounded value set, so the monodromy theorem continues
$Z$ holomorphically through $|\xi|<\sqrt S$.  It is injective there, because
$Z(\xi_1)=Z(\xi_2)$ forces $\xi_1^2=\xi_2^2$, and the two inverse sheets can
meet only at an excluded branch point.  Writing $Z(\xi)=\sum_{k\ge1}a_k\xi^k$
and fixing $1<R<\sqrt S$, injectivity and the area formula give
\[
  \pi\sum_{k\ge1}k|a_k|^2R^{2k}=\operatorname{Area}Z(D_R),
\]
while $f(Z(\xi))=1-\xi^2$ confines the image to $\{|f|<1+R^2\}$, whose area
is bounded by $\pi(1+R^2)^{2/n}$ for a monic $f$.  That bound is the
composition of two published inequalities: Dubinin's Theorem~1
\cite[Thm.~1]{dubinin}, which bounds the area of $\{|f|<s\}$ by $s^{2/n}$ times
the area of $\{|f|<1\}$, and P\'olya's classical inequality
\cite{polya1928}, which bounds the latter by $\pi$.  Cauchy--Schwarz then gives
\[
  \sum_{k\ge1}|a_k|
    \le(1+R^2)^{1/n}\Bigl(\sum_{k\ge1}\frac1{kR^{2k}}\Bigr)^{1/2}
    =(1+R^2)^{1/n}\sqrt{\log\frac{R^2}{R^2-1}},
\]
and termwise integration of $|Z'|$ over $[-1,1]$ bounds the connector length by
$2\sum_{k\ge1}|a_k|$.  Letting $R\uparrow\sqrt S$ proves the theorem.

\paragraph{Boundary.}\label{bdry:critical-value-separation}
The analytic continuation, the univalence, the area formula, and the
area--capacity inequality are ordinary mathematics; they are proved in
\path{ErdosProblems/Erdos1041/FirstMergeCriticalValueSeparationCertificate.md}
and are not formalised.  What Lean checks is the numerical half of the
statement, in
\path{ErdosProblems/Erdos1041/FirstMergeCriticalValueSeparation.lean}:
the squared coefficient
$\pdecl{firstMergeSquaredCoefficient}$ is antitone in the degree
(\pdecl{firstMergeSquaredCoefficient\_antitone\_degree}), a squared connector
length below $4$ times a coefficient below $1$ forces length below $2$
(\pdecl{firstMerge\_length\_lt\_two\_of\_squared\_bound}), and the three
regimes of Corollary~\ref{res:critical-value-thresholds} hold exactly
(\pdecl{firstMerge\_exact\_convenient\_thresholds}).  The Comparator entry is
\path{ExternalVerification1041FirstMergeCriticalValueSeparation}.  Lean does
not prove the analytic hypothesis $\ell^2\le4\,\mathrm{c}(n,S)$ that those
consumers take as input, and nothing in this section asserts the unrestricted
problem or the complementary near-tie regime.

This regime is incomparable with Pendyala's degree-four theorem
\cite[Thm.~1, p.~1]{june2026}, which is unconditional on the critical spectrum
and confined to $n=4$.  Degree four with a near-tie in the critical spectrum
is his and not ours; degree forty with $S=4$ separation is ours and not his.

\section{The Newton value equation}\label{sec:newton}

Away from the critical set, define the complex Newton field
\[
  N(z)=-\frac{f(z)}{f'(z)} .
\]
Let $z(t)$ be differentiable with $z'(t)=N(z(t))$, and put $w(t)=f(z(t))$.

\begin{theorem}[value equation]\label{res:value}
$w'(t)=-w(t)$.
\end{theorem}

\noindent The computation is one line: $w'=f'(z)\,z'=f'(z)\cdot(-f(z)/f'(z))=-f(z)=-w$.
The kernel checks it as \pdecl{newtonFlow_value_hasDerivAt}, together with the
differential form of the first integral,
\pdecl{newtonFlow_scaledValue_hasDerivAt_zero}:
\[
  \frac{d}{dt}\Bigl(e^{t}f(z(t))\Bigr)=0,
  \qquad\text{equivalently}\qquad
  f(z(t))=e^{-t}f(z(0)) .
\]
Observe what this says about the geometry.  The value moves radially inward at
exponential rate and never changes argument.  The lemniscate $\{|f|<1\}$ is
therefore forward-invariant, and the flow lines are exactly the preimages of
rays from the origin.

\begin{corollary}[ray separation]\label{res:ray}
The endpoints of any finite Newton-flow connection lie on the same oriented ray
from $0$.  Consequently, if two critical values lie on distinct positive rays,
no Newton-flow saddle connection joins the corresponding critical points.
\end{corollary}

\noindent This is checked in consumer form: the kernel accepts the implication
from the hypothesis of distinct rays to the absence of a connection.  It is the
statement the topology needs, and it is a genuine sharpening of the criterion
used in the literature.

\section{Arguments, not moduli}\label{sec:arguments}

It is tempting to arrange a generic perturbation so that the critical values
are pairwise distinct, or that their moduli are pairwise distinct, and to
conclude that saddle connections are excluded.  Neither is enough.

Two distinct critical values can lie on one ray, and two critical values with
distinct moduli certainly can: the ray records the argument, and the modulus is
exactly the coordinate the flow contracts.  By Corollary~\ref{res:ray} the
invariant that excludes connections is the argument.  What a perturbation must
therefore achieve is pairwise distinct critical-value \emph{arguments}, which is
a condition on $n-1$ points modulo the circle rather than on their positions in
the plane.

The cost of that condition is also checked, and it is small.

\begin{theorem}[ray-collision locus]\label{res:locus}
Let $a\ne b$ be complex.  Every common translation $\beta$ for which $a+\beta$
and $b+\beta$ lie on the same positive ray has the form
\[
  \beta=\frac{ra-b}{1-r},
  \qquad r\in\mathbb{R}_{>0},\ r\ne1 .
\]
\end{theorem}

\noindent So each pair of critical values contributes a one-real-parameter
forbidden locus in the translation plane, given in closed form
(\pdecl{translated_samePositiveRay_parameterization}).  A finite union of such
loci has empty interior, which is the shape one wants for an avoidance
argument.  Turning that into a perturbation of $f$ requires root retention and
length slack.  Section~\ref{sec:open} settles finite-exceptional coefficient
genericity and quantitative root retention; global component stability and
the transfer of strict metric slack remain open.

\section{Exact reciprocal coordinates and zero-margin transfer}
\label{sec:reciprocal}

The local near-polygon analysis becomes substantially cleaner when its
nonlinear terms are retained as coordinates.  Assume throughout this section
that
\[
 f(z)=\prod_{j=1}^n(z-a_j),\qquad 0<|a_j|<1,
\]
and put
\[
 r=\min_j|a_j|,\qquad p_m=\sum_{j=1}^n a_j^{-m}\quad(m\ge1).
\]

\begin{theorem}[exact reciprocal Newton expansion with $C^2$ tails]
\label{res:reciprocal-expansion}
For $|z|<r$,
\[
 \log|f(z)|
 =\log|f(0)|-\Re\sum_{m\ge1}\frac{p_m}{m}z^m,          \tag{R1}
\]
and the series is absolutely convergent.  Write $q=|z|/r<1$ and
\[
 T_N(z)=\sum_{m>N}\frac{p_m}{m}z^m.
\]
Then
\begin{align*}
 |T_N(z)|
 &\le \frac{nq^{N+1}}{(N+1)(1-q)},                     \tag{R2}\\
 |T_N'(z)|
 &\le \frac nr\frac{q^N}{1-q},                         \tag{R3}\\
 |T_N''(z)|
 &\le \frac n{r^2}q^{N-1}
       \left(\frac N{1-q}+\frac q{(1-q)^2}\right).     \tag{R4}
\end{align*}
The derivative bounds are read for $N\ge1$; the value bound holds for
$N\ge0$.
\end{theorem}

\begin{proof}
Factorwise real logarithms avoid any global choice of a complex logarithm.
The factorisation
\[
 \frac{f(z)}{f(0)}=\prod_j\left(1-\frac z{a_j}\right)
\]
and the absolutely convergent scalar series give
\[
 \log|f(z)|-\log|f(0)|
 =\sum_j\log\left|1-\frac z{a_j}\right|
 =-\Re\sum_j\sum_{m\ge1}\frac1m
       \left(\frac z{a_j}\right)^m
 =-\Re\sum_{m\ge1}\frac{p_m}{m}z^m.
\]
This proves (R1).  Since $|p_m|\le nr^{-m}$, summing the geometric majorant
gives (R2).  Termwise differentiation and the identities for the first two
differentiated geometric tails give (R3)--(R4).
\end{proof}

The $C^2$ endpoint is load-bearing: moving quadratic contacts are invisible
to a mere uniform-value estimate.  Moreover, the coordinates $p_m$ do not
depend on a labelling choice and are rational in the polynomial
coefficients.

\begin{proposition}[coefficient rationality and the exact near-Fekete bridge]
\label{res:reciprocal-rationality}
Write $f(z)=\sum_{k=0}^n c_kz^k$, $c_n=1$, $c_0\ne0$, and
$\widehat c_k=c_k/c_0$, extended by zero for $k>n$.  Then
$\widehat c_0=1$ and, for every $m\ge1$,
\[
 p_m=-m\widehat c_m-
      \sum_{j=1}^{m-1}\widehat c_jp_{m-j}.              \tag{R5}
\]
Thus every reciprocal coordinate is a rational function of the coefficients
of $f$, with denominator a power of $f(0)$.

Now write a near-regular $n$-gon as
\[
 a_k=\omega^k(1+e_k),\qquad
 \omega=e^{2\pi i/n},\qquad \eta=\max_k|e_k|<1,
\]
and put $E_m^{(j)}=\sum_ke_k^j\omega^{-km}$ and
$E_m=E_m^{(1)}$.  For $1\le m\le n-1$,
\[
 -\frac{p_m}{m}
 =E_m-\frac{m+1}{2}E_m^{(2)}
   +\frac{(m+1)(m+2)}6E_m^{(3)}-\cdots,                \tag{R6}
\]
whose $j$th term is
\[
 (-1)^{j-1}\frac1m\binom{m+j-1}{j}E_m^{(j)}.
\]
In particular,
\[
 \left|-\frac{p_m}{m}-E_m\right|
 \le \frac{n(m+1)\eta^2}{2(1-\eta)^{m+2}}.             \tag{R7}
\]
\end{proposition}

\begin{proof}[Proof spine]
Newton's identity for the roots $a_j^{-1}$ of the reversed polynomial gives
(R5) after using
$e_k(a^{-1})=(-1)^k\widehat c_k$.  For (R6), expand
$(1+e_k)^{-m}$ by the binomial series.  The constant term cancels because
$\sum_k\omega^{-km}=0$ for $1\le m<n$; the remaining terms are exactly those
displayed.  Taylor's formula with integral remainder for
$(1+e)^{-m}$ gives (R7).
\end{proof}

The conceptual point is stronger than the $O(\eta^2)$ estimate: the
``anchored nonlinear remainder'' is the second and higher part of the exact
coefficient $-p_m/m$.  It is not discarded error.  This matters precisely at
a tied Newton face, where a quadratic correction can have the same weight as
a nominally active linear mode.  The reciprocal chart retains both.

The coefficient-rational coordinates also give a compact local scale, but the
normalisation is not itself rational or semialgebraic: it includes the real
logarithm below and a local root branch.  Near the regular $n$-gon, on a
neighbourhood where $p_n\ne0$, choose a continuous nonzero branch
$\lambda^n=p_n/n$, put
\[
 b_m=(p_m/m)\lambda^{-m},\qquad
 d=-\log|f(0)|,
\]
and set
\[
 \rho=\max\left\{d^{1/n},
   \max_{1\le m<n}|b_m|^{1/(n-m)}\right\},
 \quad A_0=d/\rho^n,
 \quad A_m=b_m/\rho^{n-m}.
\]
Then $0\le A_0\le1$, $|A_m|\le1$, and at least one coordinate lies on the
boundary.  After $z=\rho\lambda^{-1}\zeta$, retaining terms through any
$N\ge n$ leaves, for each fixed $M$ and all sufficiently small $\rho$ with
$\rho M/(r|\lambda|)<1$, a $C^2$ error of order
$O_{n,N,M}(\rho^{N+1-n})$ on $|\zeta|\le M$.  Hence the unmodelled tail can
be pushed beyond any prescribed finite contact order.  The chart does not
itself prove that every limiting face transfers.  In the present open-disc
setting $d>0$ and hence $\rho>0$; the closure point $\rho=0$, corresponding
to the regular polygon after normalisation, is handled separately rather than
by dividing by the scale.

Two exact elementary consequences close the local zero-margin bookkeeping.

\begin{proposition}[staple identity and contact discriminant]
\label{res:staple-contact}
For distinct roots $a_i,a_j$, put $u_k=a_k/|a_k|$.  If
$0\le s\le\min(|a_i|,|a_j|)$, the equal-radius staple
\[
 [a_i,su_i]+[su_i,su_j]+[su_j,a_j]
\]
has length
\[
 L=|a_i|+|a_j|-s\bigl(2-|u_i-u_j|\bigr)
 \le |a_i|+|a_j|<2.                                    \tag{R8}
\]
The strict inequality holds for every admissible $s$, including $s=0$,
whenever the roots lie in the open unit disc.

If $\kappa>0$ and
\[
 F(x)\le-\kappa x^2+B|x|+V+\varepsilon,
\]
then
\[
 F(x)\le V+\varepsilon+
   \frac{\max\{B,0\}^2}{4\kappa}.                    \tag{R9}
\]
Consequently the contact remains strictly negative whenever
$V+\varepsilon+\max\{B,0\}^2/(4\kappa)<0$.
\end{proposition}

\begin{proof}
The two radial arms of the staple have lengths $|a_i|-s$ and $|a_j|-s$,
while its chord has length $s|u_i-u_j|$, proving (R8); use
$|u_i-u_j|\le2$.  Completing the square in
$-\kappa x^2+B|x|$ gives the exact supremum
$\max\{B,0\}^2/(4\kappa)$ and proves (R9).  When $B\ge0$ this is the familiar
$B^2/(4\kappa)$ formula.
\end{proof}

Proposition~\ref{res:staple-contact} rules out two recurrent false moves:
root-direction changes need not spend a polynomial-independent length margin,
and uniform convergence to a nonpositive quadratic-contact model is not
enough.  The local transverse jet and curvature must satisfy the displayed
sufficient negativity test.  Because $F$ is only bounded above by the
quadratic majorant, this test is not asserted to be necessary for contact
survival.

Theorem~\ref{res:reciprocal-expansion} and
Propositions~\ref{res:reciprocal-rationality}--\ref{res:staple-contact} are
ordinary exact theorems.  Lean checks the reciprocal power-sum majorant, the
geometric tail, the solved Newton recursion, the staple equality and strict
length consequence, and the completed-square bound and negativity consumer.
It does not formalise the complex logarithmic series, the full $C^2$ tail,
the near-Fekete bridge, the compact normalisation, or the general stratified
transfer theorem.  The classical expansion and Newton identity carry no
novelty claim; the paper contribution here is the exact identification (R6)
and the resulting coordinate choice at the transfer frontier.

\section{A proof gap in the unrestricted argument}\label{sec:gap}

The March manuscript's Proposition~12 claims the following load-bearing
statement.  For $u=-\log|f|$, a connected component $V$ of $\{u>c\}$ carrying
$m\ge2$ simple zeros, and the stated regularity and Morse hypotheses, it
constructs, for every $\varepsilon>0$, an embedded spanning tree
$G_\varepsilon\subset V$ with
\[
 \operatorname{len}(G_\varepsilon)
 \le\frac1{2\pi}\int_{2\alpha}^{\infty}P_V(t)\,dt+\varepsilon.
\tag{5.1}\label{eq:prop12-bound}
\]
The final theorem uses this estimate, so the issue below cannot be bypassed by
calling the proposition auxiliary.  Proposition~7 of the unattributed local
revision repeats the same estimate.

The estimate itself is false.  Take the Cassini polynomial
$f_a(z)=z^2-a^2$ at $a=9/10$.  Its component contains the two roots at distance
$9/5$ and satisfies the relevant Morse and critical-value hypotheses.  Direct
level-length calculation gives the upper majorant
\[
  4\bigl(\sqrt{a^2+a}-a\bigr)<\frac{41}{25}
\]
for the right side of~\eqref{eq:prop12-bound} before $\varepsilon$, whereas
every connected set containing both roots has length at least $9/5$.  Since
$9/5-41/25=4/25$, choosing a smaller positive $\varepsilon$ contradicts the
assertion.  Thus neither a different local saddle neighbourhood nor a perfect
topological decomposition can recover the printed coefficient $1/(2\pi)$.
The exact inequality and its abstract tree-budget contradiction are checked in
\texttt{ErdosProblems.Erdos1041.CassiniTreeBudget}; the level-length majorant is
an explicit analytic input, not a kernel-checked integral evaluation.

At an interior index-one critical point $p$, the proof invokes a Morse chart
\[
 u=\mu+x^2-y^2
\]
and replaces the saddle by a three-ended neighbourhood having one connected
lower cross-section and two connected upper cross-sections.  That local model
is false as written.  Because $p\in V$ and $V$ is open, a sufficiently small
closed disc around $p$ lies entirely in $V$.  In that disc the full Morse chart
has four sectors: two components of $u>\mu$ and two components of $u<\mu$.
A global component argument cannot delete one local sector from a disc already
contained in $V$.

This independently diagnoses a proof step, but the Cassini witness above also
refutes the proposition's statement.  A different route might
cut an adjoining regular annulus along a separatrix or regular flow arc before
forming the block, retain a four-pronged saddle neighbourhood and change the
assembly, or replace the local construction by the ray-cut decomposition
proposed below.  But no repair can retain~\eqref{eq:prop12-bound}; it must pay a
positive attachment cost, select only one short pair instead of spanning every
root, or use a different global metric inequality.
The shorter descriptions of the same three-ended block do not repair the
four-sector topology.

Corollary~\ref{res:ray} supplies one independent input for a different route:
distinct critical-value arguments exclude saddle-to-saddle Newton
connections.  It does not itself prove the compact planar decomposition,
classify all orbit endpoints or provide the metric gluing estimate.  Those are
separate problems below.

\section{Finite evidence}\label{sec:finite}

A search was run over random monic polynomials with roots in the unit disc.  For
each sample the region $\{|f|<1\}$ was rasterised and shortest grid paths were
computed between every pair of roots.  The best upper bounds obtained were
\[
  \begin{array}{lrr}
    \text{degree } 5, & 500 \text{ trials}: & 1.1052648928\\
    \text{degree } 6, & 1500 \text{ trials}: & 0.8450414343\\
    \text{degree } 8, & 1500 \text{ trials}: & 0.6203916714\\
    \text{degree } 10, & 1500 \text{ trials}: & 0.4303640486 .
  \end{array}
\]
No counterexample candidate was found, and the measured values sit well below
the threshold $2$.  These are grid distances for the sampled configurations:
upper bounds on those samples, not bounds over the family.  The apparent
decrease with degree is a property of the sample, and we draw no conjecture from
it.
Future searches should target named failure modes: near-degenerate saddles,
thin necks, boundary-critical configurations and almost-connected
separatrices---and report those diagnostics.  Raster paths remain candidate
finders, never continuous certificates.

\section{Exact symmetric benchmark: complementary chords for every binomial degree}
\label{sec:binomial}

% paper-assimilation: binomial_lemniscate_complementary_chords

There is a complete all-degree model family which isolates the metric geometry
without pretending to select paths for arbitrary polynomials.  Let
$f(z)=z^n-a$, where $n\ge2$ and $0<|a|<1$.  After rotation write $a=r^n>0$,
put $c=\cos(\pi/n)$, and set
\[
 r_*=(1+c^n)^{-1/n},\qquad
 \varepsilon=(1-r^n)^{1/n}.
\]

\begin{theorem}[complementary binomial chords]
Two adjacent zeros of $z^n-a$ can be joined by an explicit polygonal path
inside $\{|z^n-a|<1\}$ of length strictly below $2$.  For $r\le r_*$ the
adjacent-root chord itself works.  For $r\ge r_*$, two radial legs and an
inner adjacent crossing chord work after an arbitrarily small radial
contraction.  These two constructions meet at $r=r_*$.
\end{theorem}

The central chord calculation is exact.  On the chord between adjacent roots,
\[
 \max |z^n-r^n|=r^n(1+c^n),
\]
with the maximum at the midpoint; hence the first construction has precisely
the threshold $r_*$.  Above it, put $t=\varepsilon/c$.  The path from $r$ to
$t$ to $t\omega$ to $r\omega$ has its crossing chord in the closed
lemniscate, and its length is
\[
 2r-2\varepsilon\tan\!\left(\frac\pi4-\frac\pi{2n}\right)<2r<2.
\]
The radius $t$ is maximal: a larger crossing chord already escapes at its
midpoint.  Contracting that chord to radius $\lambda t$ with $0<\lambda<1$
makes the containment strict while keeping the two radial connections and the
strict length bound.  For $n=2$, the adjacent-root diameter is already a
strict path of length $2r<2$.

The proof uses the sharp inequality
\[
  1+\cos(n\theta)\le2\cos^n\theta
  \qquad (|\theta|\le\pi/n),
\]
followed by elementary chord geometry and the disk inequality that justifies
the contraction.  It is an authored ordinary all-degree proof with bounded
numerical replay, not a Lean theorem or a priority claim.  Its role is a
symmetric benchmark: the exact transition says why a direct outer chord fails
and why the complementary inner chord repairs it.  It gives no path
optimality, no theorem for arbitrary monic polynomials, and no solution of
Erd\H{o}s Problem~\ref{res:problem}.

% paper-assimilation: erdos1041_sharp_collinear_root_diameter
\section{The sharp collinear theorem: order buys an exact Chebyshev extremal}
\label{sec:collinear}

Collinearity supplies the global ordering that arbitrary complex root sets
lack, and it yields a sharp all-degree solved subcase.  Put
\[
 C_n=\frac{1}{2^{n-1}\cos^n(\pi/(2n))}.
\]

\begin{theorem}[sharp collinear root-diameter bound]
\label{res:sharp-collinear-root-diameter}
Let $f$ be monic of degree $n\ge2$, with collinear zero occurrences of diameter
$D$.  Then some two zero occurrences are joined by their straight segment,
of length at most $D$, on which
\[
 |f|\le C_n\left(\frac D2\right)^n.
\]
The constant is sharp: equality is attained by the affine images of the
scaled Chebyshev root configuration with extreme roots at distance $D$.
\end{theorem}

\begin{proof}
Repeated zeros give the constant path, so suppose the zeros are distinct.
Translation, rotation, and scaling by $R=D/2$ reduce to a monic real-rooted
polynomial $q$ with roots
\[
 -1=y_1<\cdots<y_n=1,
 \qquad |f|=R^n|q|\ \text{on the root line}.
\]
Set $r_n=\cos(\pi/(2n))$ and compare $q$ with the monic polynomial
\[
 q_*(x)=\frac{T_n(r_nx)}{2^{n-1}r_n^n}.
\]
It vanishes at $\pm1$ and has $|q_*|\le C_n$ on $[-1,1]$.  Choose a point
$c_i$ of maximum $|q|$ in each adjacent gap.  If every gap maximum exceeded
$C_n$, then $h=q-q_*$ would have the alternating signs of $q(c_i)$ at the
$c_i$, hence a zero between each consecutive pair, as well as zeros at
$-1$ and $1$.  These are $n$ distinct zeros, whereas the leading terms cancel
and $\deg h\le n-1$, a contradiction.  Thus one selected gap has
$|q|\le C_n$ throughout, and scaling back proves the bound.

For sharpness take the roots of $q_*$, namely
\[
 \frac{\cos((2k-1)\pi/(2n))}{r_n},\qquad 1\le k\le n.
\]
Their extreme roots are $\pm1$, and every adjacent gap reaches $C_n$ at a
Chebyshev extremum.
\end{proof}

If $D<2$, the displayed level is below one and this segment lies in the open
lemniscate.  For $n\ge3$, the same strict level holds already at $D=2$;
degree two is the exact boundary case.  The proof is a complete authored
ordinary argument.  Two load-bearing alternation and Chebyshev-comparator
kernels are Lean checked, while the affine normalization and transport remain
ordinary; bounded replay is regression evidence only.  This proves neither
that every adjacent gap is safe nor any non-collinear analogue: the real order
used in the alternation count is precisely what the general problem lacks.

\section{An auxiliary critical-value theorem through degree five}
\label{sec:freepoint}

A second route isolates the algebraic critical-value estimate from the metric
path construction.  For $c_1,\ldots,c_m\in\overline{\mathbb D}$ put
\[
 A_j=\prod_{k=1}^m |1-\overline{c_j}c_k|,
 \qquad
 S_m=\sum_{j=1}^m A_j^{1/m}.
\]
Write $\mathrm{FP}_m$ for the assertion $S_m\le m$ for every such tuple.

\begin{theorem}[four-point free-point inequality]\label{res:fp4}
$\mathrm{FP}_4$ holds.  More precisely,
\[
 \sum_{j=1}^4
   \left(\prod_{k=1}^4|1-\overline{c_j}c_k|\right)^{1/4}
 \le4,
\]
with equality if and only if $c_1=c_2=c_3=c_4=0$.
\end{theorem}

The proof divides at $\max_j|c_j|^2=21/25$.  In the central region, the
logarithmic identities $\sum_j h_j=-E$ and
$\sum_j h_j^2\le E\,\operatorname{mean}_j L_j$ combine with an exact
Bernstein-certified row envelope and rational exponential bounds.  In the
outer region, a $K_4$ matching H\"older inequality reduces the estimate to a
one-variable radical inequality with positive exact squared slack.  The two
parts meet at the stated rational radius; neither is a numerical extrapolation
from the other.

The central argument has an all-degree form that is stronger than the uniform
radius estimate used above.  Put
\[
 L_j=-\log(1-|c_j|^2),\qquad
 M_j=\frac1m\sum_{k=1}^m-\log(1-|c_j||c_k|),
\]
and set $C(0)=1/2$ and
\[
 C(t)=\frac{e^t-1-t}{t^2}\quad(t>0).
\]

\begin{theorem}[adaptive all-degree central region]
\label{res:fp-central-radius}
For every $m\ge1$, the configuration-dependent condition
\[
 \frac1m\sum_{j=1}^m C(M_j)L_j\le1
 \tag{\textsc{rowcert}}
\]
implies $S_m\le m$, with equality only when all $c_j=0$.  In particular,
let $L_*>0$ be the unique positive solution of
\[
 e^{L_*}=1+2L_*,
\]
and put
\[
 \rho_*=(1-e^{-L_*})^{1/2}=0.8457729381\ldots .
\]
Then $S_m\le m$ for every $m$ whenever
$\max_j|c_j|\le\rho_*$.
\end{theorem}

\begin{proof}[Proof spine]
Write $q_r=\sum_kc_k^r$ and
\[
 h_j=\log A_j^{1/m},\qquad
 E=\frac1m\sum_{r\ge1}\frac{|q_r|^2}{r}.
\]
The absolutely convergent logarithmic expansion gives
\[
 \sum_jh_j=-E,\qquad
 h_j^2\le\frac EmL_j,\qquad |h_j|\le M_j.
\]
On $[-M,M]$ the sharp quadratic remainder bound is
\[
 e^h\le1+h+C(M)h^2.
\]
Summing this row by row yields
\[
 S_m\le m-E\left(1-\frac1m\sum_jC(M_j)L_j\right)\le m.
\]
Equality forces every $h_j$ and hence every power sum $q_r$ to vanish;
Newton's identities then force all $c_j=0$.  If
$R=\max_j|c_j|$ and $L=-\log(1-R^2)$, then
$M_j,L_j\le L$ and $C(L)L\le1$ is exactly
$e^L\le1+2L$.  This proves the uniform-radius consequence.
\end{proof}

Thus any counterexample to $\mathrm{FP}_m$ must leave the explicit ball
$\max_j|c_j|\le\rho_*$ and must fail the sharper rowwise certificate.  This
is an unconditional infinite-degree region, but it is not a proof of
$\mathrm{FP}_m$ without restrictions for $m\ge5$.

The free-point statement enters the polynomial problem through a separate
all-degree implication.  Let $(S)_n$ denote the following assertion: if $f$ is
monic of degree $n$, the minimum enclosing radius of its roots is $R$, and
$c_1,\ldots,c_{n-1}$ are its critical points counted with multiplicity, then
\[
  \sum_{j=1}^{n-1}|f(c_j)|^{1/n}\le(n-1)R.
\]

\begin{theorem}[free-point to critical values]\label{res:fp-to-s}
For every $n\ge2$, $\mathrm{FP}_{n-1}$ implies $(S)_n$.
\end{theorem}

\begin{proof}[Proof spine]
After translating and scaling the minimum enclosing root disk, regard the
left-hand side as a symmetric function of the roots in the closed unit
polydisc.  The critical values give a continuous plurisubharmonic multiset
sum: away from critical-point collisions this follows from local holomorphic
branches, and local boundedness gives the extension across the collision
locus.  Applying the one-variable maximum principle in each root coordinate
moves the maximum to the distinguished torus.

When all roots lie on the unit circle, the critical points satisfy the exact
identity
\[
 |f(c_j)|=\prod_{k=1}^{n-1}|1-\overline{c_j}c_k|.
\]
Put $m=n-1$ and $y_j=|f(c_j)|^{1/m}$.  The hypothesis gives
$\sum_jy_j\le m$, while the tangent inequality
\[
 y^{m/(m+1)}\le\frac{my+1}{m+1}\qquad(y\ge0)
\]
converts the exponent and yields $\sum_j|f(c_j)|^{1/n}\le m$.  Scaling back
restores the factor $R$.
\end{proof}

The previously proved $\mathrm{FP}_3$ case and
Theorem~\ref{res:fp4} therefore establish $(S)_n$ for $n=2,3,4,5$; the new
cases are degrees four and five.  The proofs are ordinary exact analytic
arguments with deterministic replay.  Lean checks the scalar matching,
stationary-point, and exponent-conversion kernels, but not the complete
central/outer proof, plurisubharmonic removability, or the torus identity.

This advances the algebraic half of the programme but supplies neither a
root-pair selection nor a path-length bound.  In particular, $(S)_n$ does not
imply Problem~\ref{res:problem}.  The next free-point case is
$\mathrm{FP}_5$.

\section{The exact nonlinear lifetime currency}
\label{sec:orlicz}

The merge-tree route has its own exact theorem.  For each nontrivial child
component $u$, let $k_u$ be its number of roots and let
$r_u=\beta_u/\beta_{\operatorname{parent}(u)}\in(0,1]$ be its merge-scale
ratio.  Its contribution to the attachment age of each descendant root is
\[
 x_u=\frac1{k_u}\log\frac1{r_u}.
\]
The edge lifetime is written
\[
 \lambda_k(q)=\log\frac{1+q^{2/k}}{1-q^{2/k}},
 \qquad
 I_k(r)=\int_r^1\frac{dq}{q\lambda_k(q)}.
\]

\begin{theorem}[exact attachment-age/lifetime currency]
\label{res:orlicz-currency}
Define
\[
 \Phi(x)=\int_0^x\frac{dt}{\log(\coth t)}\qquad(x\ge0).
\]
At $t=0$ the integrand is understood by its continuous limiting value $0$
(equivalently, the integral is improper at that endpoint).
For every integer $k\ge1$ and $0<r\le1$, with
$x=k^{-1}\log(1/r)$,
\[
 I_k(r)=k\Phi(x).                                      \tag{O1}
\]
The function $\Phi$ is increasing and strictly convex on $(0,\infty)$, and
\[
 \frac{\Phi(x)}x\longrightarrow0\qquad(x\downarrow0). \tag{O2}
\]
Consequently, for every fixed $k\ge1$ and every $c>0$, some $0<r<1$
satisfies
\[
 I_k(r)<c\,\frac1k\log\frac1r.                        \tag{O3}
\]
In particular, there is no positive universal linear conversion from
attachment age to lifetime.
\end{theorem}

\begin{proof}
Set $q=e^{-kt}$.  Then $dq/q=-k\,dt$, the endpoints $q=1,r$ become
$t=0,x$, and
\[
 \frac{1+e^{-2t}}{1-e^{-2t}}=\coth t,
\]
which proves (O1).  Since $\coth t$, and hence $\log(\coth t)$, is strictly
decreasing, the positive integrand $1/\log(\coth t)$ is strictly increasing.
This proves monotonicity and strict convexity.  Moreover
\[
 0\le\frac{\Phi(x)}x\le\frac1{\log(\coth x)}\longrightarrow0,
\]
which gives (O2); taking $r=e^{-kx}$ gives (O3).
\end{proof}

The same monotonicity gives a useful non-linear lower bound that the no-linear
floor does not erase.  For $x>0$,
\[
 \Phi(x)\ge\int_{x/2}^x\frac{dt}{\log(\coth(x/2))}
 =\frac{x}{2\log(\coth(x/2))},
\]
and therefore
\[
 I_k(r)\ge
 \frac{kx}{2\log(\coth(x/2))}.                       \tag{O3$'$}
\]
The point is distribution-sensitive curvature, not a uniform positive linear
coefficient at the origin.

This identity retains the complete distribution of edge ages.  If a root
chain has edges $u=1,\ldots,m$, set
\[
 A=\sum_u x_u,\qquad K=\sum_u k_u,\qquad
 L=\sum_u I_{k_u}(r_u).
\]
Then the exact chain ledger and weighted Jensen inequality give
\[
L=\sum_u k_u\Phi(x_u)
 \ge K\Phi\!\left(\frac{\sum_u k_ux_u}{K}\right)
 \ge K\Phi\!\left(\frac{2A}{K}\right),                \tag{O4}
\]
where the last inequality uses $k_u\ge2$ for ancestor edges.  Separately,
unweighted Jensen and the same inequality $k_u\ge2$ give the
cardinality form
\[
 L\ge2\sum_u\Phi(x_u)\ge2m\Phi(A/m).                  \tag{O4$'$}
\]
It is not obtained by weakening the weighted $K$-bound: it records different
data.  For a strict component chain below an $n$-root top one also has the
exact combinatorial ceilings
\[
 m\le n-2,
 \qquad
 K\le2+3+\cdots+(n-1)=\frac{n(n-1)}2-1.
\]
These estimates are
degree-dependent by design: splitting a fixed age among many very rapid
mergers is exactly the regime in which a linear estimate loses all force.

The stronger global fact comes from summing across roots instead of following
one root at a time.  Fix an internal component $v$ with root set $S_v$, and
let $A_i(v)$ be the sum of the increments $x_u$ along the path from root $i$
to $v$.  Let $u\preceq v$ range over the nontrivial proper descendant
components whose parent still lies below or equals $v$.

\begin{theorem}[root-summed identity and two-young-root selector]
\label{res:orlicz-selector}
One has the exact identity
\[
 \sum_{i\in S_v}A_i(v)
 =\sum_{u\preceq v}\log\frac1{r_u}
 =\log\prod_{u\preceq v}\frac1{r_u}.                  \tag{O5}
\]
If $|S_v|\ge2$, there are distinct $i,j\in S_v$ such that
\[
 \max\{A_i(v),A_j(v)\}
 \le
 \frac{\sum_{u\preceq v}\log(1/r_u)}{|S_v|-1}.       \tag{O6}
\]
The constant $1/(|S_v|-1)$ is sharp for the underlying ordered-age
inequality.
\end{theorem}

\begin{proof}
The edge $u\to\operatorname{parent}(u)$ contributes
$k_u^{-1}\log(1/r_u)$ to exactly the $k_u$ roots below it.  Its coefficient
after summing over roots is therefore one, proving the first equality in
(O5); the second is the logarithm product law.  Order the nonnegative ages as
$a_1\le a_2\le\cdots\le a_{|S_v|}$.  The largest $|S_v|-1$ entries are at
least $a_2$, so
$(|S_v|-1)a_2\le\sum_i a_i$.  Choosing the first two roots proves (O6).
For arbitrary nonnegative age data the limiting vector $(0,c,\ldots,c)$
shows sharpness of this ordered-age step.
\end{proof}

The theorem is an unconditional attachment-compatible pair selector and an
exact diagnosis of the metric frontier.  It does \emph{not} upper-bound the
internal product $\prod_{u\preceq v}1/r_u$, show that the selected roots are
compatible with one short embedded path, or perform the saddle and annular
gluing.  Those are the actual remaining producers.  The source note proves
the substitution, Jensen, and double-counting arguments and supplies bounded
deterministic replay; Lean checks only the lifetime-to-Orlicz division, the
weighted-chain order transfer, the abstract no-linear-floor implication, and
the final second-age division.

\section{Complements and further questions}\label{sec:open}

The dependency chain has five separate gates.  A proof, a minimally corrected
hypothesis set, or an explicit polynomial or planar counterexample is a
decisive answer to any one of them.

\subsection*{1. Repair or refute the saddle block}

\begin{problem}[corrected local saddle assembly]\label{prob:saddle1041}
Under the hypotheses of Proposition~12, replace the invalid three-ended local
model by a four-pronged block or by a block formed after a specified annular
cut.  Prove that for every $\eta>0$ its connector has total Euclidean length
below $\eta$, uniformly over the selected attachment points, and that the
corrected blocks assemble to an embedded tree satisfying
\[
 \operatorname{len}(G_\varepsilon)
 \le\frac1{2\pi}\int_{2\alpha}^{\infty}P_V(t)\,dt+\varepsilon;
\]
or give a polynomial or harmonic planar counterexample to that statement.
\end{problem}

The Cassini family in \S\ref{sec:gap} answers this problem negatively as
stated: the requested global estimate is smaller than the distance between its
two roots.  What remains meaningful is a weaker one-pair assembly with a
different metric budget; changing only the local saddle block cannot restore
the coefficient in the displayed inequality.

The full-disc model $x^2-y^2$, an annulus in which lower branches rejoin, two
saddles joined by a separatrix and simultaneous saddle levels are mandatory
tests.  Repeating the one-lower/two-upper assertion does not answer the
problem.

\subsection*{2. The compact ray-cut decomposition}

Let $p$ have simple roots, simple nonzero critical points and let
$u=-\log|p|$.  For regular values $c<T$, take
\[
 M_{c,T}=\overline{V\cap\{c<u<T\}},
\]
and assume explicitly that this is a compact genus-zero surface, its lower
boundary is one smooth Jordan curve, its upper boundary consists of $m$ smooth
root curves, all interior critical points are nondegenerate index-one saddles,
the normalised gradient field
\[
 X=\frac{\nabla u}{|\nabla u|^2}=-\frac{p}{p'}
\]
is transverse to the level boundaries, and no maximal $X$-trajectory has two
saddle endpoints.

\begin{problem}[finite strip decomposition after ray cuts]\label{prob:reeb1041}
For every $\eta>0$, construct disjoint Morse neighbourhoods $N_j$ with
$\sum_j\operatorname{diam}N_j<\eta$ and a finite set of complete separatrix or
regular-flow cuts so that every component of the complement is flow-diffeomorphic
to a rectangle
\[
 [a_S,b_S]\times[0,1],\qquad u(\Phi_S(t,s))=t,
\]
with connected level sections and no uncut annular component.  Give an explicit
finite bound $E(s,m)$ for the number of strips, or exhibit the minimal missing
hypothesis or a counterexample.
\end{problem}

No particular formula such as $2s+1$ is presumed.  Boundary tangencies,
simultaneous levels, branch reunion through an annulus and non-Hausdorff orbit
spaces must be handled rather than suppressed.

\subsection*{3. Metric fan-in without losing the coefficient}

For a strip $S$, write
\[
 \Gamma_t^S=S\cap\{u=t\},\qquad
 P_S(t)=\mathcal H^1(\Gamma_t^S),
\]
and let $k_S$ be its number of root ends.  The flux identity
\[
 \int_{\Gamma_t^S}|\nabla u|\,ds=2\pi k_S
\]
gives the average trajectory estimate
\[
 \int_{\Gamma_{t_0}^S}\operatorname{len}(\gamma_x)\,d\mu_{t_0}(x)
 =
 \frac1{2\pi k_S}\int_{a_S}^{b_S}P_S(t)\,dt.
\]

\begin{problem}[additive-error strip gluing]\label{prob:metric1041}
Assuming Problem~\ref{prob:reeb1041}, select trajectories in all strips and
connect them through the saddle and root neighbourhoods so that, for every
$\eta>0$, the resulting embedded tree contains all $m$ roots and obeys
\[
 \operatorname{len}(G)
 \le\frac1{2\pi}\int_{2\alpha}^{\infty}P_V(t)\,dt+\eta.
\tag{6.1}\label{eq:metric-fanin1041}
\]
The total saddle, annular-cut and root-cap cost must be below $\eta$ without a
multiplicative loss in $1/(2\pi)$.
\end{problem}

For the final strict inequality, use the actual collar slack
\[
 q=\frac1{2\pi}\int_{\alpha}^{2\alpha}P_V(t)\,dt>0
\]
and give budgets that keep the perturbation, tree error and transfer cost below
fixed fractions of $q$.  Independently choosing a shortest trajectory in each
strip is not enough.  Theorems~\ref{res:orlicz-currency} and
\ref{res:orlicz-selector} identify the exact missing control: a linear
lifetime charge is impossible, while an upper bound on one internal
merge-ratio product would select two young roots before the attachment age
fragments across many edges.

\subsection*{4. Coefficient perturbation and stability}

The constant-translation stage is no longer open.  Once a finite critical-value
family is injective, Lean proves an arbitrarily small translation making every
value nonzero and pairwise positive-ray separated
(\pdecl{exists_small_translation_separating_arguments}).  It also proves the
explicit root-retention estimate; a shift below $\varepsilon$ keeps all roots
in the unit disc when
\[
 ((n+1)\varepsilon)^{1/n}+\rho<1
\]
(\pdecl{constant_perturbation_roots_in_unitDisk}).  A constant translation
cannot separate initially equal critical values.

\begin{proposition}[a one-parameter critical-value genericity lemma]
\label{prop:linear-genericity1041}
Let $f\in\mathbb C[z]$ have degree at least $2$.  There is a finite set
$E\subset\mathbb C$ such that, for every $\lambda\notin E$, the polynomial
$f+\lambda z$ has only simple critical points and its critical values are
pairwise distinct.  In particular, such $\lambda$ occur arbitrarily close to
$0$.
\end{proposition}

\begin{proof}
Write $d=\deg f$ and let $a$ be the leading coefficient.  A multiple critical
point of $f+\lambda z$ satisfies
\[
 f'(u)+\lambda=0,\qquad f''(u)=0.
\]
Since $f''$ is not the zero polynomial, this excludes only finitely many
values $\lambda=-f'(u)$.

For distinct critical points $u,v$ with equal critical values, eliminating
$\lambda$ gives the two polynomial equations
\[
 A(u,v):=\frac{f'(u)-f'(v)}{u-v}=0,
 \qquad
 B(u,v):=\frac{f(u)-f(v)}{u-v}-f'(u)=0,
\tag{4.1}\label{eq:critical-collision}
\]
and then $\lambda=-f'(u)$.  These two equations have only finitely many
solutions.  Indeed, their top homogeneous parts are
\[
 A_{d-2}=da\,(u^{d-2}+u^{d-3}v+\cdots+v^{d-2}),
\]
and
\[
 B_{d-1}=a\,(u^{d-1}+u^{d-2}v+\cdots+v^{d-1}-d u^{d-1}).
\]
They have no common projective zero at infinity.  For $u\ne0$, putting
$r=v/u$, the first equation gives
$1+r+\cdots+r^{d-2}=0$, hence $r^{d-1}=1$ and $r\ne1$; the second expression,
divided by $a u^{d-1}$, is then $1-d\ne0$.  If $u=0$, the first expression is
nonzero (and for $d=2$ it is already a nonzero constant).  Thus $A$ and $B$
have no common curve component: any such affine component would have a point
at infinity.  Their common zero set is finite, so the associated values
$-f'(u)$ form a finite exceptional set.  Avoiding its union with the first
finite set proves the claim.
\end{proof}

\begin{proposition}[compact barrier stability]\label{prop:compact-barrier1041}
Let $C,H\subset\mathbb C$ be compact, let $D_1,\ldots,D_m$ be pairwise
disjoint closed discs, and let $T\subset\mathbb R_{>0}$ be finite.  Suppose
that $C$ and $H$ are nonempty, that $m\geq1$, that $T$ is nonempty,
that $f'$ has no zero on $C$, that $f$ has no zero on any $\partial D_i$, and
that $|f(z)|\ne t$ for every $z\in H$ and $t\in T$.  Put
\[
 c_C=\min_C|f'|,\qquad
 c_D=\min_i\min_{\partial D_i}|f|,\qquad
 c_H=\min_{t\in T}\min_H\bigl||f|-t\bigr|,
\]
and
\[
 R=\max\bigl\{|z|:z\in H\cup\partial D_1\cup\cdots\cup\partial D_m\bigr\}.
\]
Then $c_C,c_D,c_H>0$.  If
\[
 |\lambda|<c_C,
 \qquad
 R|\lambda|+|\beta|<\min(c_D,c_H),
\tag{4.2}\label{eq:compact-barrier-budget}
\]
then $g_{\lambda,\beta}=f+\lambda z+\beta$ has no critical point in $C$,
has the same number of zeros as $f$ in every $D_i$ (with multiplicity), and
lies on the same side of every level $t\in T$ as $f$ at every point of $H$.
\end{proposition}

\begin{proof}
Compactness and the hypotheses give the three positive minima.  Since
$g'_{\lambda,\beta}=f'+\lambda$, the first inequality in
\eqref{eq:compact-barrier-budget} gives
\[
 |g'_{\lambda,\beta}(z)|\ge |f'(z)|-|\lambda|>0
 \qquad(z\in C).
\]
On every $\partial D_i$ the second inequality gives
\[
 |\lambda z+\beta|\le R|\lambda|+|\beta|<|f(z)|,
\]
so Rouch\'e's theorem preserves the zero count in $D_i$.  Finally,
\[
 \bigl||g_{\lambda,\beta}(z)|-|f(z)|\bigr|
 \le |\lambda z+\beta|<\bigl||f(z)|-t\bigr|
 \qquad(z\in H,\ t\in T),
\]
which preserves the sign of $|f(z)|-t$.
\end{proof}

\noindent This compact estimate separates several local requirements that
were previously bundled into ``stability'': protected collars and truncation
boundaries remain free of critical points, prescribed root discs retain their
zeros, and finite level barriers retain their chosen side.  It does not assert
that a selected sublevel component remains connected or nested, and it gives
no continuity statement for the coarea slack $q$.
The first, protected-set noncriticality implication is separately Lean-checked
in \texttt{noncritical\_on\_of\_norm\_lt\_uniform\_lower\_bound} and is
targeted by the external Comparator interface.  The Rouch\'e root-count and
level-side implications in this proposition remain ordinary paper proofs.

\begin{problem}[two-stage generic perturbation with remaining global stability]\label{prob:perturb1041}
For fixed root discs, compact collar $K$, regular levels
$\alpha/2,\alpha,2\alpha,5\alpha/2$ and collar slack $q>0$, prove that an
arbitrarily small
\[
 g_{\lambda,\beta}(z)=f(z)+\lambda z+\beta
\]
can be chosen so that:

\begin{enumerate}
\item $f+\lambda z$ has simple relevant critical points and injective complex
critical values;
\item the checked constant translation $\beta$ makes those values nonzero and
pairwise ray-separated;
\item the relevant component remains in $K$, contains exactly the corresponding
perturbed roots and has slack $q_g>q/2$; and
\item root displacement and straight-line transfer back to the original roots
consume less than a prescribed fraction of $q/m$.
\end{enumerate}

Proposition~\ref{prop:linear-genericity1041} proves the first item, with an
arbitrarily small one-coefficient perturbation.  Proposition~\ref{prop:compact-barrier1041}
supplies the local collar, truncation, root-disc and finite-level controls once
their compact margins have been fixed.  The remaining issue is global component
and strict-slack stability, followed by the quantified root-transfer budget.
\end{problem}

Finite planar avoidance and the subsequent constant translation must not be
relisted as missing.  Critical-value genericity and compact barrier control are
now elementary inputs; global component stability, strict slack stability and
the transfer budget remain the open content.

\subsection*{5. The global Newton-flow claim ceiling}

\paragraph{What the constant-factor theorem leaves open.}
Theorems~\ref{res:constant-factor-path} and
\ref{res:first-merge-capacity-closures} remove the former absence of any
unconditional global path bound.  They also close substantial parent regimes.
They do not prove the sharp constant $2$ on the remaining low-arity,
high-component-capacity cells or on the high-critical branch.  The Newton-flow
programme below is therefore no longer needed to obtain a finite global
constant; its purpose is to recover sharp containment and sharp metric
allocation on precisely those residual cells.  P\'olya's primary theorem
supplies only the global area--capacity input used in
Section~\ref{sec:constant-factor}; no priority claim is inferred for the
source-specific local estimates or assembled paths.

\begin{problem}[relative global Newton-flow theorem]\label{prob:globalflow1041}
On a compact regular band
\[
 M_{a,b}=\{z:a\le-\log|p(z)|\le b\},
\]
prove real-time existence and the identities
\[
 p(z(t))=e^{-(t-t_0)}p(z(t_0)),\qquad
 u(z(t))=u(z(t_0))+t-t_0
\]
throughout every maximal orbit; classify both limiting endpoints among the
regular boundary, root ends and finite saddle set; and exclude recurrent,
accumulating and non-Hausdorff orbit phenomena.  Under pairwise ray separation
of the critical values, decide whether the flow is Morse--Smale relative to the
boundary or whether the strongest valid conclusion is only absence of
saddle-to-saddle connections.
\end{problem}

The checked algebra proves the pointwise value equation and the endpoint-ray
consumer.  It does not supply global solution theory or an orbit-space graph.
A positive stronger theorem must give the graph and a finite edge bound; a
negative answer should exhibit the simplest ray-separated polynomial carrying
the remaining pathology.

The auxiliary free-point route now has a different boundary.  Theorem
\ref{res:fp4} closes $\mathrm{FP}_4$,
Theorem~\ref{res:fp-central-radius} proves an explicit all-degree central
region and a sharper adaptive certificate, and
Theorem~\ref{res:fp-to-s} transfers $\mathrm{FP}_{n-1}$ to $(S)_n$ in every
degree.  What remains there is unrestricted $\mathrm{FP}_m$ for $m\ge5$ and,
independently, a metric or selection theorem that converts critical-value
control into the short root-to-root path.

Erd\H{o}s~\#1041 remains open.  The source now publicly verifies the Newton
kernel, finite ray avoidance, quantitative constant-translation root control,
and the protected-set noncriticality inequality in the compact-barrier
argument; ordinary proofs supply finite-exceptional coefficient genericity and
the remaining root-disc and level-barrier controls.  Global component
stability, the corrected planar decomposition and metric gluing remain the
exact unresolved producers.

\section{Evidence and boundary ledger}
\label{sec:evidence-ledger}

This ledger separates proved statements, checked kernels, computations, and
open steps.  Each entry states both the evidence and the limit of the claim.

{\small
\renewcommand{\arraystretch}{1.16}
\begin{longtable}{@{}
  >{\raggedright\arraybackslash}p{0.28\linewidth}
  >{\raggedright\arraybackslash}p{0.66\linewidth}@{}}
\toprule
\textbf{Result or mechanism} & \textbf{Evidence and exact boundary} \\
\midrule
\endfirsthead
\toprule
\textbf{Result or mechanism} & \textbf{Evidence and exact boundary} \\
\midrule
\endhead
\midrule
\multicolumn{2}{r@{}}{\footnotesize\itshape Continued on the next page} \\
\endfoot
\bottomrule
\endlastfoot
\multicolumn{2}{@{}p{0.94\linewidth}@{}}{\textbf{Principal results}} \\
\addlinespace
Critical two-root proximity &
\textsc{Lean checked.}  At every non-root critical point, two distinct roots
have total critical-point distance at most twice the geometric mean of all
root distances.  The theorem selects a metric pair; it does not prove that a
short connecting path stays in the lemniscate. \\
\addlinespace
Straight-path obstructions &
\textsc{Lean-checked counterexamples.}  Exact families rule out both the
unique-nearest spoke and every pair of straight spokes through the critical
point as universal constructions. \\
\addlinespace
Translated cubic quotient fibres &
\textsc{Checked cubic kernel; ordinary finite-fibre assembly.}  This gives a
complete solved family $f(z)=P((z-h)^q)$.  Lean proves the safe cubic spoke;
the root-of-unity assembly is an ordinary proof. \\
\addlinespace
Primitive sparse quintics &
\textsc{Checked selector kernels; ordinary path assembly.}  Two tail-energy
selectors give the solved open-disc family.  The moment-to-path transport is
an ordinary proof. \\
\addlinespace
Sharp collinear family &
\textsc{Checked kernels; ordinary affine transport.}  This is a complete
all-degree family with the sharp Chebyshev constant. \\
\midrule
\multicolumn{2}{@{}p{0.94\linewidth}@{}}{\textbf{Critical-value and free-point route}} \\
\addlinespace
Newton identities &
\textsc{Lean checked.}  Away from critical points, a trajectory tangent to
$-f/f'$ satisfies $w'=-w$ and
$\frac{d}{dt}(e^t f(z(t)))=0$. \\
\addlinespace
Critical-value rays &
\textsc{Lean checked, consumer form.}  Endpoints of a finite connection share
one oriented ray, so distinct rays exclude a connection.  For a pair of
critical values, collision occurs on
$\beta=(ra-b)/(1-r)$ with $r>0$ and $r\ne1$. \\
\addlinespace
Quartic case &
\textsc{Cited theorem.}  The degree-four result is proved in
\cite[Thm.~1, p.~1]{june2026}.  It does not extend to arbitrary degree. \\
\addlinespace
$\mathrm{FP}_4$ &
\textsc{Ordinary exact proof; deterministic replay.}  Four free points in
$\overline{\mathbb D}$ satisfy the product-distance mean inequality, with
equality only at the all-zero tuple.  Lean checks scalar certificate kernels,
not the full analytic proof. \\
\addlinespace
All-degree central region &
\textsc{Ordinary exact proof; deterministic replay.}  The adaptive rowwise
certificate proves the free-point inequality in every degree; uniformly it
covers $\max_j|c_j|\le0.8457729381\ldots$.  This restricts every possible
higher-degree counterexample but does not prove unrestricted
$\mathrm{FP}_m$ for $m\ge5$. \\
\addlinespace
$\mathrm{FP}_{n-1}\Rightarrow(S)_n$ &
\textsc{Ordinary exact proof; deterministic replay.}  The proof uses a
plurisubharmonic reduction to the root torus, the torus product identity, and
an exact exponent conversion.  The several-complex-variables step is not
formalised in Lean. \\
\addlinespace
Critical-value inequality $(S)_n$ &
\textsc{Proved for $2\le n\le5$.}  Degrees four and five are new here.  This
auxiliary inequality alone does not produce a root-to-root path. \\
\midrule
\multicolumn{2}{@{}p{0.94\linewidth}@{}}{\textbf{Attachment-time and local-transfer route}} \\
\addlinespace
Attachment-age currency &
\textsc{Ordinary exact proof; checked scalar consumers.}
$I_k(r)=k\Phi(k^{-1}\log(1/r))$, and weighted Jensen bounds hold.  No
positive uniform linear age charge can hold. \\
\addlinespace
Two-young-root selector &
\textsc{Ordinary exact proof; finite regression.}  Subtree-size denominators
cancel exactly, and two roots meet the sharp total-age threshold.  The result
does not supply a merge-product or geometric compatibility bound. \\
\addlinespace
Reciprocal Newton expansion &
\textsc{Ordinary exact proof; checked finite kernels.}  The expansion has
coefficient-rational Taylor coordinates and explicit $C^0$, $C^1$, and $C^2$
tails.  It retains the anchored nonlinear remainder coefficientwise. \\
\addlinespace
Staple and quadratic contact &
\textsc{Ordinary exact proof; checked scalar kernels.}  The staple spends no
length margin.  Completing the square gives a sharp quadratic-majorant bound
and a sufficient negativity criterion for the underlying contact.  A global
stratified transfer theorem remains open. \\
\addlinespace
Global constant-factor path &
\textsc{Ordinary exact proof; finite regression only.}  Two zero occurrences
admit a path of length at most $(71/10)\mu^{1/n}$ in $K_{2\mu}$; when
$\mu\le1/2$, the length is at most $5.7$ in $\{|f|<1\}$.  This result has no
Lean or independent-review claim. \\
\addlinespace
First-merge parent regimes &
\textsc{Ordinary exact proof.}  The target constant $2$ holds at
$\mu\le1/2$ for first-merge arity at least $17$, and under the explicit
arity-dependent component-capacity cutoffs of
Theorem~\ref{res:first-merge-capacity-closures}.  The complementary cells and
the high-critical branch remain open. \\
\midrule
\multicolumn{2}{@{}p{0.94\linewidth}@{}}{\textbf{Genericity, failed mechanisms, and open steps}} \\
\addlinespace
March 2026 proof mechanism &
\textsc{Refuted.}  A Cassini witness refutes the tree budget in
\cite[Prop.~12]{march2026}; the local saddle block also fails independently. \\
\addlinespace
Constant translation &
\textsc{Lean checked.}  Critical-value injectivity permits arbitrarily small
ray avoidance while retaining an explicit unit-disc root margin. \\
\addlinespace
Coefficient perturbation &
\textsc{Partly formalised.}  The note proves finite-exceptional genericity;
Lean checks protected-set noncriticality.  Root-disc and level-barrier control
remain ordinary arguments, while global component and strict-slack stability
remain open. \\
\addlinespace
Reeb decomposition and length fan-in &
\textsc{Open.}  These are the two remaining producers in the Newton-flow
route. \\
\addlinespace
Random search through degree $10$ &
\textsc{Verified finite instances.}  The computations give upper bounds only
for the sampled configurations. \\
\addlinespace
Erd\H{o}s \#1041 &
\textsc{Open.}  No proof of the general conjecture is claimed. \\
\end{longtable}}

\section*{Statements and declarations}

This manuscript is authored exposition, not Lean proof authority.  Its
strongest unrestricted-geometry paper theorem is the ordinary analytic
constant-factor path theorem, together with the first-merge arity and
component-capacity regimes that reach the parent constant $2$.  Lean does not
formalise their coarea, univalent-map, component-capacity, or path-assembly
arguments.

Its
strongest source-current checked family is the critical geometric-mean
proximity theorem
\pdecl{exists_two_roots_dist_sum_le_two_mul_geomMean}, paired with the exact
unique-nearest-spoke and all-straight-cubic no-go certificates.  The translated
cubic quotient-fibre theorem, primitive sparse-quintic theorem, and sharp
collinear theorem are complete ordinary paper theorems whose load-bearing
selectors or alternation kernels are Lean checked.  Lean does not formalise
their finite-fibre, moment-to-path, or affine-transport assemblies.

The older checked core is the Newton value equation, the exponential first
integral, the consumer form of ray separation, the finite planar-avoidance
theorem, quantitative constant-translation root retention, and the
ray-collision parameterisation.  The decomposition and length statements of
\S\ref{sec:open} are not proved.  The
diagnosis of Proposition~12 in \S\ref{sec:gap} has two independent parts: the
printed local saddle construction is topologically invalid, and the exact
Cassini budget refutes the proposition's metric statement.  The
search results of \S\ref{sec:finite} are computations.

The four-point free-point inequality, the adaptive all-degree central-region
theorem, and the all-degree $\mathrm{FP}_{n-1}\Rightarrow(S)_n$ bridge are
ordinary exact analytic theorems with green deterministic replay.  The files
\texttt{FreePointFP4Complete.lean} and
\texttt{FreePointTorusPshReduction.lean} check scalar algebraic kernels only.
They do not formalise the complete central/outer analysis, plurisubharmonic
removability, the torus product identity, or a root-to-root path.  No novelty
or literature-priority claim is made for these auxiliary inequalities.

The exact attachment-age/lifetime transform, its Jensen consequences, the
root-summed age/product identity, and the two-young-root selector are ordinary
exact theorems proved in the source note
\texttt{AttachmentAgeLifetimeOrlicz.md}.  The companion Lean module checks
only four scalar consumers after the analytic identities and monotonicity have
been supplied.  The five Comparator packages expose those four consumer
types: \texttt{OrliczNoUniformLinear} and \texttt{NoUniformLinear} are
duplicate interfaces for one mathematical obstruction, not two distinct
paper results.  No claim is made here that Lean formalises the integral
substitution, convexity, Jensen's inequality, the merge-tree double count, or
the remaining geometric allocation theorem.

The reciprocal Newton expansion, coefficient recursion, near-Fekete bridge,
compact normalisation, staple identity, and completed-square contact criterion
are ordinary exact theorems proved in
\texttt{ReciprocalNewtonExpansion.md}.  Lean checks seven distinct finite
kernel types exposed through nine Comparator packages:
\texttt{ReciprocalTail} duplicates \texttt{NormRecipPowerSum}, and
\texttt{ContactAbsorption} duplicates \texttt{ContactSurvives}.  These are
one reciprocal-coordinate and transfer family, not nine independent paper
claims.  The logarithmic-series identity, full $C^2$ bounds, exact
near-Fekete identification, and rational compactification remain ordinary
proof rather than Lean-formalised assertions.

\appendix
\section{Guide to the formal sources}\label{app:sources}

The source-current checked summit is distributed across
\texttt{CriticalTwoRootProximity}, \texttt{CubicQuotientFiberCase},
\texttt{PrimitiveQuinticBoundaryTail},
\texttt{PrimitiveQuinticInteriorTail},
\texttt{SharpCollinearAlternation}, and
\texttt{SharpCollinearChebyshev}.  The public
\texttt{ErdosProblems.Erdos1041.NewtonFlowRaySeparation} module contains the
older checked Newton-flow source.  The search of
\S\ref{sec:finite} is
\texttt{scripts/search\_counterexample.py} in the source package.  The
declaration table below is pinned to the shared formal-source commit used for
the Newton-flow declarations in this problem-note series.

\begin{itemize}[leftmargin=*]
\item \lref{Erdos1041/CriticalTwoRootProximity.lean}{289}{exists_two_roots_dist_sum_le_two_mul_geomMean}
\item \lref{Erdos1041/CriticalTwoRootProximity.lean}{438}{nearestSpoke_unique_nearest_spoke_escapes}
\item \lref{Erdos1041/CriticalTwoRootProximity.lean}{569}{allStraightCubic_every_pair_midpoint_escapes}
\item \lref{Erdos1041/CubicQuotientFiberCase.lean}{161}{cubic_has_safe_root_spoke}
\item \lref{Erdos1041/PrimitiveQuinticBoundaryTail.lean}{200}{primitiveBoundary_exists_two_tailEnergy_lt_one}
\item \lref{Erdos1041/PrimitiveQuinticInteriorTail.lean}{272}{primitiveInterior_exists_two_tailEnergy_lt_one}
\item \lref{Erdos1041/SharpCollinearAlternation.lean}{158}{exists_peak_le_of_monic_comparison}
\item \lref{Erdos1041/SharpCollinearChebyshev.lean}{133}{exists_peak_le_comparisonBound}
\item \lref{Erdos1041/NewtonFlowRaySeparation.lean}{34}{newtonFlowVector}
\item \lref{Erdos1041/NewtonFlowRaySeparation.lean}{38}{derivative_mul_newtonFlowVector}
\item \lref{Erdos1041/NewtonFlowRaySeparation.lean}{50}{newtonFlow_value_hasDerivAt}
\item \lref{Erdos1041/NewtonFlowRaySeparation.lean}{64}{newtonFlow_scaledValue_hasDerivAt_zero}
\item \lref{Erdos1041/NewtonFlowRaySeparation.lean}{77}{SamePositiveRay}
\item \lref{Erdos1041/NewtonFlowRaySeparation.lean}{80}{samePositiveRay_refl}
\item \lref{Erdos1041/NewtonFlowRaySeparation.lean}{84}{samePositiveRay_symm}
\item \lref{Erdos1041/NewtonFlowRaySeparation.lean}{92}{samePositiveRay_trans}
\item \lref{Erdos1041/NewtonFlowRaySeparation.lean}{107}{translated_samePositiveRay_parameterization}
\item \lref{Erdos1041/NewtonFlowRaySeparation.lean}{127}{realAffineLine}
\item \lref{Erdos1041/NewtonFlowRaySeparation.lean}{130}{realAffineLine_eq_affineSpan}
\item \lref{Erdos1041/NewtonFlowRaySeparation.lean}{147}{isClosed_realAffineLine}
\item \lref{Erdos1041/NewtonFlowRaySeparation.lean}{152}{dense_compl_realAffineLine}
\item \lref{Erdos1041/NewtonFlowRaySeparation.lean}{162}{exists_small_avoiding_finite_realAffineLines}
\item \lref{Erdos1041/NewtonFlowRaySeparation.lean}{179}{samePositiveRay_imp_mem_realAffineLine}
\item \lref{Erdos1041/NewtonFlowRaySeparation.lean}{197}{exists_small_translation_separating_arguments}
\item \lref{Erdos1041/NewtonFlowRaySeparation.lean}{230}{norm_lt_one_of_near_root}
\item \lref{Erdos1041/NewtonFlowRaySeparation.lean}{257}{constant_perturbation_root_near_original}
\item \lref{Erdos1041/NewtonFlowRaySeparation.lean}{287}{constant_perturbation_roots_in_unitDisk}
\item \lref{Erdos1041/NewtonFlowRaySeparation.lean}{306}{samePositiveRay_of_real_exp_decay}
\item \lref{Erdos1041/NewtonFlowRaySeparation.lean}{315}{no_newtonConnection_of_not_samePositiveRay}
\end{itemize}

\begingroup
\raggedright
\paragraph{Source-current auxiliary proofs.}
The proofs of Theorems~\ref{res:fp4} and
\ref{res:fp-central-radius} are
\path{ErdosProblems/Erdos1041/FreePointZeroInsertionFP4Matching.md}; its
deterministic replay is
\path{scripts/check_erdos1041_free_point_zero_insertion_fp4.py}.
The scalar Lean companion contains
\path{three_mul_pairProduct_le_sum_sq} and the five
\path{outer_*} stationary-point certificates.  The proof of
Theorem~\ref{res:fp-to-s} is
\path{ErdosProblems/Erdos1041/FreePointTorusPshReduction.md}; its replay is
\path{scripts/check_erdos1041_free_point_torus_psh_reduction.py}.
Its Lean companion checks
\path{rpow_le_tangent_at_one},
\path{sum_rpow_le_card_of_sum_le_card}, and the two finite exponent
bridge declarations.  These source-current files are not covered by the older
Newton-flow commit pin below; the paper states their separate authority and
validation boundary explicitly.

The attachment-age/lifetime proof is
\path{ErdosProblems/Erdos1041/AttachmentAgeLifetimeOrlicz.md}, and its scalar
Lean companion is
\path{ErdosProblems/Erdos1041/AttachmentAgeLifetimeOrlicz.lean}.  Its
Comparator interfaces are
\path{ExternalVerification1041LifetimeOrliczBound},
\path{ExternalVerification1041WeightedChainConsumer},
\path{ExternalVerification1041SecondAgeBound}, and
\path{ExternalVerification1041OrliczNoUniformLinear}.  These interfaces cover the
formalised consumers, not the ordinary analytic core of
Theorems~\ref{res:orlicz-currency}--\ref{res:orlicz-selector}.

The reciprocal-coordinate proof and replay description are
\path{ErdosProblems/Erdos1041/ReciprocalNewtonExpansion.md}; its finite Lean
kernel is
\path{ErdosProblems/Erdos1041/ReciprocalNewtonExpansion.lean}, and its exact
replay is
\path{ErdosProblems/Erdos1041/scripts/check_erdos1041_reciprocal_newton_expansion.py}.
The Comparator interfaces are
\path{ExternalVerification1041NewtonSolve},
\path{ExternalVerification1041NormRecipPowerSum},
\path{ExternalVerification1041GeometricTail},
\path{ExternalVerification1041QuadraticContact},
\path{ExternalVerification1041ContactSurvives},
\path{ExternalVerification1041StapleEquality}, and
\path{ExternalVerification1041StapleLength}.  They cover the seven checked
finite kernels identified above, not every ordinary theorem in
Section~\ref{sec:reciprocal}.

The finite near-Fekete result described in the abstract is the ordinary
analytic source note
\texttt{ErdosProblems/Erdos1041/NearFeketeTransverseClosure.md}, pinned to
source commit
\texttt{0f248d8826eca5f946052c1d4fd2f983bc3d7a04} and SHA-256
{\footnotesize\nolinkurl{3fbfcab1310c5d596990baab102d5787331f1df37a85b8becb2aab77eb5daadd}}.
Its limited formal consumer is
\texttt{ErdosProblems.Erdos1041.NearFeketeTransverseClosure}, whose tracked
Lean source is pinned to commit
\texttt{ee2bfdebfe6d7d4827cf0d1b243f37fd3667c3e7} and SHA-256
{\footnotesize\nolinkurl{20cd39f660e19c71bb47868721e7f69b6ae5585045087237b0bd6a9fa64d43a5}}.
The declarations at lines 19--60 check only eventual sign and length transfer;
they do not formalize the two-cutoff complex analysis, tangent strata, bulk
regime, or the unrestricted problem.

\paragraph{Immutable source-generation boundary.}
The declaration table above is pinned to formal-source commit
\texttt{7e80158d80c9beccbaa930a0b45da2db70f253d4} (23 August 2026), whose
tracked source file
\texttt{ErdosProblems/Erdos1041/NewtonFlowRaySeparation.lean} has SHA-256
{\footnotesize\nolinkurl{b3cde5cb40fb93fbbd410e8931cd9ae6fcb4e30e1346ac62d62ebaee55f65fb1}}.
This pin covers only the declarations named above, including the exact
value-equation, ray-separation, finite-avoidance, and root-retention locators;
it is not a claim that the surrounding global-flow or length argument is
formalised.  The bounded project entrypoint for replay is
\texttt{formal\_math/erdos257\_period\_noncollapse/scripts/lean\_fast\_build.py}
with the module target
\texttt{ErdosProblems.Erdos1041.NewtonFlowRaySeparation}; a replay of that
target remains validation evidence for the listed declarations, not an
independent proof of the paper's ordinary arguments.
\par
\endgroup

\begin{thebibliography}{10}
\bibitem{bloom} T.~F. Bloom, \emph{Erd\H{o}s Problems}, problem~1041.
  \url{https://www.erdosproblems.com/1041}
\bibitem{ehp1958} P.~Erd\H{o}s, F.~Herzog, and G.~Piranian,
  \emph{Metric properties of polynomials}, J. Analyse Math. \textbf{6}
  (1958), 125--148. \url{https://doi.org/10.1007/BF02790232}
\bibitem{march2026} \texttt{shtuka},
  \emph{A Short Path Joining Two Zeros Inside a Polynomial Lemniscate},
  manuscript posted 24 March 2026, 48~pp.
  \url{https://shtuka123.github.io/1041/main.pdf}
\bibitem{june2026} V.~S. Pendyala,
  \emph{A Degree-Four Lemniscate Path Theorem},
  arXiv:2606.24875v1 (2026).
  \url{https://arxiv.org/abs/2606.24875},
  doi:\href{https://doi.org/10.48550/arXiv.2606.24875}
  {10.48550/arXiv.2606.24875}.
\bibitem{pendyala2026shortest} V.~S. Pendyala,
  \emph{Shortest paths in polynomial lemniscate sublevel sets and a problem
  of Erd\H{o}s}, arXiv:2606.19178v1 (2026),
  Theorem~1.2 and introductory discussion, pp.~1--2.
  \url{https://arxiv.org/abs/2606.19178},
  doi:\href{https://doi.org/10.48550/arXiv.2606.19178}
  {10.48550/arXiv.2606.19178}.
\bibitem{dubinin} V.~N. Dubinin, \emph{Some inequalities for polynomials and
  rational functions associated with lemniscates}, Zap. Nauchn. Sem. POMI
  \textbf{404} (2012), 83--99; English translation, J. Math. Sci. \textbf{193}
  (2013), no.~1, 45--54,
  doi:\href{https://doi.org/10.1007/s10958-013-1432-4}{10.1007/s10958-013-1432-4}.
  Theorem~1 is the P\'olya-type area inequality: the area of the sublevel
  region $\{|f|<s\}$ is at most $s^{2/n}$ times the area of $\{|f|<1\}$.  It
  was identified by T.~Tao on the Erd\H{o}s Problem~\#1041 discussion page,
  25 March 2026, as the source of that scaling factor.  The published text is
  paywalled and has been seen only in Tao's paraphrase, so the inequality is
  used here exactly in the form just stated and no sharper consequence is
  drawn from it.
\bibitem{polya1928} G. P\'olya,
  \emph{Beitrag zur Verallgemeinerung des Verzerrungssatzes auf mehrfach
  zusammenh\"angende Gebiete}, Sitzungsberichte der Preussischen Akademie der
  Wissenschaften, Physikalisch-Mathematische Klasse (1928), printed
  pp.~228--232 and 280--282.
  \url{https://archive.org/details/sitzungsbericht1928preu}.
\end{thebibliography}

\end{document}
